% !TEX program = lualatex
\documentclass[11pt,a4paper]{ltjsarticle}
\usepackage{amsmath,amssymb,amsthm,mathtools}
\usepackage{luatexja-fontspec}
\usepackage[colorlinks=true,urlcolor=blue,linkcolor=blue]{hyperref}
\setmainjfont{HaranoAjiMincho-Regular}
\newtheorem{theorem}{定理}[section]
\newtheorem{lemma}[theorem]{補題}
\newtheorem{proposition}[theorem]{命題}
\newtheorem{corollary}[theorem]{系}
\newtheorem{remark}[theorem]{注意}
\renewcommand{\proofname}{証明}
\newcommand{\CC}{\mathbb C}
\newcommand{\Spec}{\operatorname{Spec}}
\newcommand{\End}{\operatorname{End}}
\newcommand{\Hom}{\operatorname{Hom}}
\newcommand{\Mat}{\operatorname{Mat}}
\newcommand{\gldim}{\operatorname{gldim}}
\newcommand{\Frac}{\operatorname{Frac}}
\newcommand{\CM}{\operatorname{CM}}
\title{中心三次商の正則性と非可換$A_1$局所模型\\
\large PBW変形、全中心、量子的判別式とBrauer障害}
\date{2026年11月13日・v8}
\begin{document}
\maketitle
\begin{abstract}
三変数Sklyanin代数の定数項変形について、PBW性と自己双対な双加群解消を示し、
三次元Calabi--Yau性を導く。指定した一混合項族の周囲の代数は
全係数値でKoszul AS正則かつNakayama twistが恒等であることを、
相対PBW性と$\mathbb Q[\alpha]$上の有限証明書で確認する。
非斉次全中心の関係式を有理数上で厳密に検証し、Rees退化でAzumaya locusを同定する。
零中心ファイバーの四つの非可換$A_1$模型と一般Brauer指数6、
二中心chartの21中心特異点と三つのホモロジー的特異点、
射影特異圏の三因子分解および球面twistを説明する。
中心chartと導分によるperfect性の貼合せから、指定混合族の
正確なtails判別式$(\alpha^3+9)(28\alpha^3+243)$を得る。
さらに分離型二変数三次式の必要十分判別式、固定した良い周囲の
超越的中心値と有限位数楕円境界での非空正則開集合を証明する。
零中心ファイバーをCayley三次曲面の反標準補集合として明示し、twisted NCCRと解釈する。
基礎となる斉次Sklyaninの性質、既知の分類・双対性と本稿の具体的な証明・計算を区別する。
\end{abstract}

% 基礎となる斉次Sklyanin代数からPBW・CY入力を導く。
\section{定数項変形のPBW性と自己双対な三次元解消}
\label{sec:quotient-pbw-cy}
以下では斉次Sklyanin代数
\[
 \begin{gathered}
 S_1=\CC\langle X,Y,Z\rangle/(r_X,r_Y,r_Z),\\
 r_X=YZ-ZY-X^2,\quad r_Y=ZX-XZ-Y^2,\quad r_Z=XY-YX-Z^2
 \end{gathered}
\]
のKoszul AS正則性（次元3、Hilbert級数$(1-t)^{-3}$）を既知入力とする。
この斉次代数の基本的な性質は\cite[Proposition~2.21, \S8.2]{WWYSklyanin}を参照する。
非斉次代数のPBW性とCY性を独立した未証明の仮定として追加する必要はない。
これらは\cite[Theorem~3.3.2]{EGDelPezzo}の既知の一般的なpotential変形定理の特殊化でもある。
以下では具体的なsyzygyと自己双対解消による証明、および被約基底上での相対化を与える。

\begin{proposition}[任意の定数項変形のPBW性]
\label{prop:quotient-constant-pbw}
任意の$c_X,c_Y,c_Z\in\CC$に対し
\[
 R(c)=\CC\langle X,Y,Z\rangle/
 (r_X-c_X,r_Y-c_Y,r_Z-c_Z)
\]
は語長filtrationについて$\operatorname{gr}R(c)=S_1$を満たす。
特に$R_1=R(1,0,0)$について$\dim F_nR_1=\binom{n+3}{3}$である。
\end{proposition}
\begin{proof}
$V=\langle X,Y,Z\rangle$、$\mathcal R=\langle r_X,r_Y,r_Z\rangle$とする。
Koszul解消から$(V\otimes\mathcal R)\cap(\mathcal R\otimes V)$は一次元で、
\[
 \eta=X\otimes r_X+Y\otimes r_Y+Z\otimes r_Z
     =r_X\otimes X+r_Y\otimes Y+r_Z\otimes Z
\]
で生成される。最後の等式は六つの三文字語と三つの立方の展開から直接従う。
Braverman--Gaitsgory\cite[Main Theorem~0.5, Lemma~3.3]{BGPBW}の記号で
一次補正は零、定数補正は$\beta(r_i)=c_i$である。
同定理の条件(i),(iii)は自動的に成立し、条件(ii)も
\[
 (\beta\otimes1-1\otimes\beta)(\eta)
 =c_XX+c_YY+c_ZZ-Xc_X-Yc_Y-Zc_Z=0
\]
から従う。従ってPBW性が成立する。
\end{proof}

\begin{proposition}[被約有限型基底上の相対PBW性]
\label{prop:quotient-relative-pbw}
$C$を被約有限生成可換$\CC$-代数、$c_X,c_Y,c_Z\in C$とする。
同じ表示で定義した$C$-代数$R_C(c)$は
\[
 \operatorname{gr}R_C(c)\simeq C\otimes_\CC S_1
\]
を満たし、斉次語基底を持ち上げると$C$上自由である。
\end{proposition}
\begin{proof}
$V_n$を長さ$n$以下の自由語の$C$加群に、
$|a|+|b|\le n-2$の関係式倍$a(r_i-c_i)b$を課した有限表示加群とする。
各閉点への特殊化は前命題のPBW代数であり、$\dim(V_n\otimes k(s))=\binom{n+3}{3}$。
ここでPBW性から、次数$n$以下の関係はこれら有界次数の倍で生成される。
実際、関係の最高次項を斉次関係の倍で消して次数を下げればよい。

$V_n$の有限表示行列のrankは全閉点で一定となる。
それより大きいminorは被約性とNullstellensatzで$C$上零であり、
最大rankのminorの非消滅開集合は$\operatorname{Spec}C$を覆う。
各開集合で行列を恒等blockと零blockに変形すると、$V_n$は有限射影となり、
形成は任意の基底変更と可換する。
$V_n\to V_{n+1}$は全閉点で単射で、rankも一定なので、同じ議論から
単射で有限射影なcokernelを持つ。

$\varinjlim V_n=R_C(c)$なので$V_n$はその実際の語長filtrationに等しい。
自然な全射$C\otimes(S_1)_n\to V_n/V_{n-1}$は
全閉点で同型であり、有限射影加群の同rankの写像なので同型である。
各商は自由だからfiltrationの短完全列は分裂し、語基底の持上げが
$C$上の自由基底を与える。
被約性を外して閉点のrankだけから同じ結論を推論してはいけない。
\end{proof}

\begin{proposition}[定数項変形の三次元Calabi--Yau性]
\label{prop:quotient-constant-cy}
$R(c)$はhomologically smoothで、
\[
 \mathbf R\operatorname{Hom}_{R(c)^e}(R(c),R(c)^e)
 \simeq R(c)[-3].
\]
従って特に$R_1$は三次元Calabi--Yauである。
\end{proposition}
\begin{proof}
$A=R(c)$とし、$P_i=A\otimes E_i\otimes A$を、
$E_0=\CC$、$E_1=V$、$E_2=\mathcal R$、$E_3=\CC\eta$で定める。
微分を
\[
 d_1(1\otimes X_i\otimes1)=X_i\otimes1-1\otimes X_i,
\]
$d_2(1\otimes r_i\otimes1)=\partial(r_i)$、
\[
 d_3(1\otimes\eta\otimes1)
 =\sum_i(X_i\otimes r_i\otimes1-1\otimes r_i\otimes X_i)
\]
とする。ここで$\partial(w)$は、語$w$の各文字を順に真ん中の$V$へ移し、
その前後の語を両端へ置く普遍微分である。
$d_1d_2=0$は$\partial(r_i)$のtelescopingと$r_i=c_i$から、
$d_2d_3=0$は$\eta$の二表示と定数の中心性から従う。
各$E_i$に次数$i$を与えたfiltrationのassociated gradedは
$S_1$のKoszul双加群解消である。
最高次項を消して次数を下げる帰納法により
\[
 0\to P_3\to P_2\to P_1\to P_0\to A\to0
\]
は完全であり、有限自由双加群解消となる。

自己双対性も直接確認できる。
$L_i=X_i\otimes1$、$R_i=1\otimes X_i^{\mathrm{op}}$、
$D_i=L_i-R_i$と書くと、$d_2$の各基底像の係数行列は
\[
 \begin{pmatrix}
 -(L_X+R_X)&D_Z&-D_Y\\
 -D_Z&-(L_Y+R_Y)&D_X\\
 D_Y&-D_X&-(L_Z+R_Z)
 \end{pmatrix}.
\]
この行列は、両端を交換する反同型$\tau(L_i)=R_i$と転置の合成で不変である。
$d_1$と$d_3$も同じ交換で符号を除き互いに双対となる。
従って上の複体の$A^e$-双対は、標準の複体の符号を調整して
元の複体を三つずらしたものとなる。
右端の余核の双加群は通常の$A$であり、Nakayama twistは生じない。
これで表示した導来双対とsmoothnessを同時に得る。
\end{proof}

\begin{corollary}[指定混合族の周囲と中心三次商]
\label{cor:quotient-ambient-pbw}
全ての$\alpha\in\CC$について
\[
 B_\alpha=\CC[u,v]\langle x,y,z\rangle/
 ([x,y]-z^2,[y,z]-x^2-\alpha uv,[z,x]-y^2)
\]
はnoetherianなKoszul AS正則整域（次元5、index 5）である。
$\Omega_1$を本文の対称化中心三次式とすると
$W_\alpha=\Omega_1+\alpha uvx+u^3+v^3$は中心正則元で、
$A_\alpha=B_\alpha/(W_\alpha)$は$\CC[u,v]$上自由であり、
全ての非零一次式$au+bv$は$A_\alpha$上正則である。
また$B_\alpha/(v-\lambda u)$は全ての$\lambda\in\CC$で
Koszul AS正則（次元4、index 4）である。
\end{corollary}
\begin{proof}
相対PBW命題を$C=\CC[u,v]$、$c_X=\alpha uv,c_Y=c_Z=0$に適用する。
$u,v$は中心正則列で$B_\alpha/(u,v)=S_1$。
相対filtrationのassociated gradedが$S_1[u,v]$なので
整域性とnoetherian性も従う。

次数1の中心正則元$z$に対するchange-of-rings Tor spectral sequenceを使う。
$B/(z)$がKoszulなら、$\operatorname{Tor}^B_q(B/(z),\CC)$は
$q=0,1$で内部次数$q$に集中するので、$E_2$は全て内部次数$p+q$に集中する。
従って$B$もKoszulで、$\operatorname{pd}_B\CC=
\operatorname{pd}_{B/(z)}\CC+1$となる。
これを$u,v$に適用して両側の射影次元5を得る。
minimal graded解消とgood filtrationで全加群の大域次元も5となる。
中心正則列に対するReesの補題からAS--Gorenstein条件とindex 5も従う。
同じ論法を基底$\CC[u]$、$c_X=\alpha\lambda u^2$に適用すれば
$B_\alpha/(v-\lambda u)$について次元4、index 4を得る。

$\Omega_1+mx$の中心性は$m=1$での厳密な中心元計算
（命題\ref{prop:quotient-verified-center6}）と、
$m\ne0$での$(x,y,z)=\sqrt m(X,Y,Z)$による規格化から従う。
$m=0$では斉次中心元である。
相対PBW基底の各係数について全閉点でこの恒等式が成立するので、
被約基底$\CC[u,v]$上でも中心的である。
$W_\alpha$は非零で、最高$x,y,z$次数部分$\Omega_1$は$S_1$上正則。
従って
\[
 \operatorname{gr}_{x,y,z}A_\alpha
 =(S_1/(\Omega_1))\otimes\CC[u,v].
\]
斉次語基底の持上げから最後の自由性と正則性を得る。
\end{proof}

\begin{proposition}[指定混合族の全係数値でのNakayama twistの消滅]
\label{prop:quotient-ambient-nakayama}
全ての$\alpha\in\CC$について$B_\alpha$のNakayama自己同型は恒等であり、
$B_\alpha$は五次元Calabi--Yau代数である。
$A_\alpha$と$C_\alpha=A_\alpha/(v-u)$のbalanced dualizing complexは
それぞれ$A_\alpha(-2)[4]$、$C_\alpha(-1)[3]$である。
従って球面twistの構成に必要なNakayama twistの恒等性も追加仮定ではない。
\end{proposition}
\begin{proof}
$E=B_\alpha^!$をKoszul双対とし、生成元を$e_x,e_y,e_z,e_u,e_v$とする。
直交関係を取ると、15本の二次関係は
\[
 e_x^2+e_ye_z=e_y^2+e_ze_x=e_z^2+e_xe_y=0,
 \qquad e_u^2=e_v^2=0,
\]
\[
 e_ie_j+e_je_i=0\quad(i<j,\ (i,j)\ne(u,v)),\qquad
 e_ue_v+e_ve_u+\alpha e_ye_z=0.
\]
$E$はHilbert級数$(1+t)^5$のgraded Frobenius代数である。
以下の有限計算は$\mathbb Q[\alpha]$上でその構造を直接検証する。
次数4の基底を
\[
 b_1=e_x^3e_u,\ b_2=e_x^3e_v,\ b_3=e_x^2e_ue_v,\
 b_4=e_xe_ye_ue_v,\ b_5=e_xe_ze_ue_v
\]
とし、次数5の基底を$h=e_x^3e_ue_v$とする。
$\lambda:E_5\to\CC$を$\lambda(h)=1$で定めると、
両方の行列$(\lambda(e_i b_j))_{i,j}$と
$(\lambda(b_j e_i))_{i,j}$は同じ
\[
 M=\begin{pmatrix}
 0&0&1&0&0\\
 0&0&0&0&1\\
 0&0&0&-1&0\\
 0&-1&0&0&0\\
 1&0&0&0&0
 \end{pmatrix},\qquad\det M=1
\]
となる。従ってFrobeniusのNakayama自己同型$\nu_E$は次数1で恒等、
次数1が代数を生成するため全体で恒等である。

再現コード\texttt{computations/verify\_central\_quotient\_koszul.py}は
次数6までの自由語商を構成し、次元$1,5,10,10,5,1,0$、
全ての相補次数のFrobenius対合の行列式$1,1,-1,-1,1,1$、
上の左右の積の一致を厳密に検査する。
全ての行簡約pivotは非零有理定数であり、$\alpha$を含む式で割らない。
従って基底と等式は全ての複素$\alpha$へ特殊化でき、一般点だけの計算ではない。
同コードは三次syzygyと前命題の自己双対な双加群複体の積も検査する。

Koszul AS正則代数の双対公式
\cite[proof of Theorem~6.3]{RRZNakayama}は
$\mu_{B_\alpha}|_{B_1}=(-1)^{5+1}(\nu_E|_{E_1})^*$を与える。
符号も正なので$\mu_{B_\alpha}=\mathrm{id}$。
\cite[Lemma~1.2]{RRZNakayama}よりAS正則性はskew CY性を与えるから、
これは通常の五次元CY性である。
中心正則元$W_\alpha$と$v-u$による商には
\cite[Lemma~1.5]{RRZNakayama}を適用でき、normalizing automorphismも恒等なので
両商のNakayama自己同型も恒等となる。
AS次元とindexは順に$(4,2)$、$(3,1)$であり、
\cite[Corollary~4.3 and its preceding discussion]{YZSerre}の
balanced dualizing complexの式を適用して結論する。
\end{proof}


% 旧中心商論文VI の第14--15節：中心元の有限表示とAzumaya locusの証明依存
% このファイルは papers/central_quotient_local_models.tex から入力する。
\section{全中心の明示生成元とRees退化によるAzumaya locus}
\label{sec:quotient-center-azumaya}
$R_1$を
\[
 \CC\langle X,Y,Z\rangle/
 ([X,Y]-Z^2,\,[Y,Z]-(X^2+1),\,[Z,X]-Y^2)
\]
とする。語長filtrationを$F_\bullet R_1$と書く。
旧中心商原稿の第14節では$\operatorname{gr}R_1=S_1$、
$\dim F_nR_1=\binom{n+3}{3}$を用いて、
次の非斉次中心元を明示する。

\begin{proposition}[四つの中心生成元：有限計算に依存する結果]
\label{prop:quotient-full-center}
$\Omega_1=(X^3+Y^3+Z^3)/3-
 (XYZ+XZY+YXZ+YZX+ZXY+ZYX)/9$と置き、
$g=\frac32(\Omega_1+X)=Y^3-ZYX+X$とする。
以下の語を定義する：
\[
 U=Z^3(g-Z^3)=Y^3(g-Y^3),
\]
\[
\begin{aligned}
 V={}&YZ^2Y^3-ZYZ^2YX-ZYZ^4
   +2Z^2Y^2ZX+Z^2YZ^3-3Z^3YZ^2\\
  &+YZ^2X+ZYZX+Z^2YX+3Z^4,\\
 W={}&YZ^4Y+ZYZ^3Y+Z^2Y^2Z^2+2Z^2YZ^2Y
      -Z^4Y^2+Z^5X\\
  &-ZYZ^2-2Z^2YZ-2Z^3Y+Z^2 .
\end{aligned}
\]
以下の有限計算を用いる：
(i) $g,U,V,W$の三生成元との交換子は$\mathbb Q$上で零、
(ii) $F_6R_1$での中心核は
$1,g,g^2,U,V,W$が基底となる六次元空間。
これらは命題\ref{prop:quotient-verified-center6}で
有理数演算と法$101$の非零rank証明を組み合わせて証明する。
(iii) 重み$18$以下の$50$中心単項式の像のrankは
法$101$では$49$で、指定の$\Phi$も法$101$で零である。
\emph{特性零における$\Phi=0$は未検証の入力}である。
これを有理数上で厳密に証明できれば、
斉次中心の既知の生成・関係定理と組み合わせて
下記の中心表示を導ける。
さらに斉次楕円Sklyanin代数$S_1$の中心構造定理
\cite{WWYSklyanin}を用いる。
この条件の下で
\[
 Z(R_1)\cong
 \frac{\CC[g,U,V,W]}{(\Phi)},\quad
 \Phi=49U^3-(13g^2+14)U^2+(g^2+1)^2U
 +(11U-g^2-1)VW+V^3-W^3
\]
であり、
$\operatorname{gr}Z(R_1)=Z(S_1)$となる。
\end{proposition}
\begin{proof}[旧原稿の証明と現行の検証範囲]
次数7以下の自由語3280個と関係式倍6015本を
有理数上で消去すると$\dim F_7R_1=120$、
$\dim F_6R_1=84$となる。
交換子写像の法$101$におけるrankは78で、
既に有理数上で中心性・独立性が確認された
上記6元から逆の上界が得られる。
従ってrank78と中心核の基底性は特性零でも確定する
（命題\ref{prop:quotient-verified-center6}）。
斉次中心定理\cite{WWYSklyanin}によれば、
$Z(S_1)$は三次の中心元と三つの六次中心元で生成され、
重み18のただ一つの関係を持つ。
$g,U,V,W$の最高次symbolがそれらを張ることを
命題\ref{prop:quotient-generators-graded}により確認した後、
任意の中心元の最高次項を順に消す次数下降帰納法により
$Z(R_1)$の生成とassociated gradedの等式を得る。

重み18以下の50単項式について、
独立の帰納的語基底アルゴリズムは
法$101$で$\dim F_{18}R_1=1330$、rank49を与える。
$\Phi$を自由語として展開し同じ商へ写した剰余も
法$101$で零と検査した。しかしこれは
\emph{有理数上の零判定の証明ではない}。
特性零における$\Phi=0$を入力として採用すれば、
法$101$のrank49から50単項式の特性零rankは少なくとも49となり、
$\Phi$による非自明な線形関係から高々49。
従って$\Phi$が唯一の重み18の関係式となる。
その最高次symbolは斉次中心の唯一の関係を生成するので、
任意の関係式の最高次項から$\Phi$の倍数を減じる
filtration帰納法により全関係イデアルが$(\Phi)$となる。
\textbf{現在の障害：} 指定した$\Phi$の有理数上の厳密恒等式を
独立に確立する必要がある。法$101$での零判定を
特性零の等式に読み替えてはいけない。
\end{proof}


\begin{proposition}[有限次数中心の独立計算証明]
\label{prop:quotient-verified-center6}
$R_1$の語長filtrationについて
\[
 \dim_{\mathbb C}F_7R_1=120,\qquad
 \dim_{\mathbb C}F_6R_1=84,\qquad
 Z(R_1)\cap F_6R_1=\langle1,g,g^2,U,V,W\rangle_{\mathbb C}.
\]
特に最後の空間の次元は6である。
\end{proposition}
\begin{proof}
次数7以下の自由語は$\sum_{i=0}^7 3^i=3280$個。
三つの定義関係式$r$の全語倍$arb$で
$|a|+|b|\le5$を取ると6015個となる。
これらが定義イデアルの$F_7$部分を生成するので、
有理数上の辞書順疎行列消去で得られる
余次元120は正確な商次元である。
同じ標準剰余から$F_6$の次元84を得る。

付属の\texttt{computations/verify\_center\_degree7.py}
は$\mathbb Q$上で$[g,x_i],[U,x_i],[V,x_i],[W,x_i]$
がすべて零となること、および
$1,g,g^2,U,V,W$が一次独立であることを
元の定義関係式から逐次計算する。
一方、この84次元空間から三生成元との交換子
\[
 F_6R_1\longrightarrow(F_7R_1)^{\oplus3}
\]
の行列を法$101$へ還元して消去するとrank78を得る。
したがって特性零でrankは少なくとも78。
上記の六つの有理数上で独立な中心元により
rankは高々$84-6=78$なので等号である。
計算コードはPython標準ライブラリの有理数演算と
法$101$の整数演算のみで再現可能である。
\end{proof}

\begin{proposition}[全中心の生成と重み18の関係の次元]
\label{prop:quotient-generators-graded}
斉次Sklyanin中心の既知の生成・関係定理
\cite[Proposition~2.31 and \S8.2]{WWYSklyanin}を入力とすると、
\[
 Z(R_1)=\CC[g,U,V,W],\qquad
 \operatorname{gr}Z(R_1)=Z(S_1).
\]
また$\deg g=3$, $\deg U=\deg V=\deg W=6$の
重み18以下の50個の中心単項式が張る空間は、
\emph{特性零でも}次元49である。
したがって重み18以下には1次元の関係空間がある。
ただし、その関係式が上の明示的な$\Phi$と
正確に一致することは、この命題の結論ではない。
\end{proposition}
\begin{proof}
非斉次関係式から定数項を消した斉次Sklyanin代数を
$S_1=\operatorname{gr}R_1$とする。
$g$の三次最高次記号と
$g^2,U,V,W$の六次最高次記号は中心元である。
後の4元が$(S_1)_6$で一次独立であることを
\texttt{computations/verify\_center\_homogeneous\_symbols.py}
の有理数上の計算で確認した。
既知の斉次中心定理では、次数三は一次元、
次数六は四次元で、中心はこれらの元によって生成される。
従って上記の記号は$Z(S_1)$の生成元を与える。

任意の$z\in Z(R_1)$の最高次記号は
$Z(S_1)$の上記生成元の多項式である。
これに対応する$g,U,V,W$の多項式を$z$から引くと、
語長filtrationの次数が下がる。これを有限回繰り返せば
$z$は$g,U,V,W$の多項式となり、
同時に$\operatorname{gr}Z(R_1)=Z(S_1)$を得る。

既知の斉次中心のHilbert級数は
\[
 \operatorname{Hilb}_{Z(S_1)}(t)
 =\frac{1-t^{18}}{(1-t^3)(1-t^6)^3}
\]
である。従って重み18以下の中心部分の次元は
自由可換代数の同範囲の50単項式から、
重み18の一つの関係を引いた49である。
上記の全中心の生成性により、これら50単項式は
$Z(R_1)\cap F_{18}R_1$を張り、
filtered associated gradedの次元保存からrankは49。
残る1次元の関係式の最高次記号は
斉次中心の関係を生成する。しかしその
\emph{非斉次の下位項}は有理数上でまだ同定していない。
\end{proof}

\begin{proposition}[Azumaya locusの同定：既知Sklyanin入力付き]
\label{prop:quotient-azumaya}
命題\ref{prop:quotient-full-center}の中心表示と
$\operatorname{gr}Z(R_1)=Z(S_1)$を仮定し、
$S_1$がPI次数6の三次元AS正則Sklyanin代数で、
その非Azumaya locusが三本の曲線からなること
\cite{WWYSklyanin}、
また$R_1$がfiltered deformationとして
左右大域次元3のprime noetherian整域であることを入力とする。
このとき
\[
 \operatorname{Az}(R_1)
 =\operatorname{Sm}(\operatorname{Spec}Z(R_1)).
\]
特に$R_1$は中心上有限な最大Cohen--Macaulay加群で、
中心上の一般rankは$36$であり、
高さ1ではAzumayaである。
\end{proposition}
\begin{proof}
$Z=Z(R_1)$、$S_1=\operatorname{gr}R_1$とし、
Rees環と中心のRees環を
\[
 \mathcal T=\bigoplus_{d\ge0}F_dR_1\,t^d,\qquad
 \mathcal C=\bigoplus_{d\ge0}(Z\cap F_dR_1)t^d
\]
と書く。元の中心のassociated gradedが
$Z(S_1)$に等しいことから、有限個の斉次中心生成元を
次数を下げつつ持ち上げられる。
従って$\mathcal T$は$\mathcal C$上有限で、
\[
 \mathcal T/(t)=S_1,\quad
 \mathcal C/(t)=Z(S_1),\quad
 \mathcal T[t^{-1}]=R_1[t,t^{-1}],\quad
 \mathcal C[t^{-1}]=Z[t,t^{-1}]
\]
が成り立つ。
$\mathcal C_{(t)}$は一次元noetherian局所整域で
極大イデアルが$t$で生成されるためDVR。
$\mathcal T_{(t)}$は有限torsion-freeゆえ自由で、
そのrankは特殊ファイバー$S_1$の中心上の一般rank
$6^2=36$に一致する。
従って$R_1$の$Z$上の一般rankも$36$。

Zariskian filtrationのgrade/GK次元比較
\cite{LevasseurNoncomm}により、
有限生成$R_1$加群$M$について
\[
 j_{R_1}(M)=j_{S_1}(\operatorname{gr}M),\quad
 \operatorname{GKdim}M=\operatorname{GKdim}\operatorname{gr}M.
\]
$S_1$はAuslander-regularかつGK-Cohen--Macaulay
なので$j_{R_1}(M)+\operatorname{GKdim}M=3$。
中心上有限であるからKrull次元とGK次元が一致し、
Brown--Macleodの基準\cite[Theorem~4.8]{BrownMacleod}
によって$R_1$は$Z$-Macaulayである。

$\mathcal C$の$(t)$を含む閉点で$S_1$がAzumayaなら、
$\mathcal T/\mathfrak m\mathcal T\simeq\operatorname{Mat}_6(\CC)$。
Nakayamaと一般rank36から
$\mathcal T_{\mathfrak m}$は$\mathcal C_{\mathfrak m}$上自由rank36で、
剰余体で同型となる
$\mathcal T_{\mathfrak m}\otimes_{\mathcal C_{\mathfrak m}}
 \mathcal T_{\mathfrak m}^{\mathrm{op}}
\to\operatorname{End}_{\mathcal C_{\mathfrak m}}(\mathcal T_{\mathfrak m})$
はNakayamaにより同型。
従って特殊中心のAzumaya点はRees族でもAzumayaである。

$\mathcal T$の非Azumaya locusを$D$とする。
次数からの$\mathbb G_m$作用で$D$は不変。
$D\cap V(t)$は$S_1$の非Azumaya locusの三曲線で次元1。
$t\ne0$と交わる$D$の既約成分の閉包は
$V(t)$にも交わり、次元公式によりその次元は高々2。
$t\ne0$で$D$は
$(\operatorname{Spec}Z\setminus\operatorname{Az}(R_1))
\times\mathbb G_m$なので、
$R_1$の非Azumaya locusは中心3次元において
次元高々1、すなわちheight one Azumaya。
Brown--Macleod\cite[Theorem~3.13]{BrownMacleod}の
prime、中心有限、有限大域次元、$Z$-Macaulay、
height one Azumayaの条件が揃い、
Azumaya locusと中心のsmooth locusが一致する。
\end{proof}

\begin{lemma}[Azumaya環と中心の大域次元]
\label{lem:quotient-azumaya-gldim}
$C$を可換noetherian環、$B$を$C$上定数正rankの
Azumaya代数とする。このとき左右大域次元は
$\operatorname{gldim}B=\operatorname{gldim}C$であり、
無限大の場合も含む。
\end{lemma}
\begin{proof}
separability idempotentにより
任意の$B$加群$M$は$B\otimes_CM$の$B$直和因子。
$B$は$C$上射影的なので
$C$-射影分解を誘導して
$\operatorname{gldim}B\le\operatorname{gldim}C$を得る。
逆向きには中心への正規化traceで$C\hookrightarrow B$が
$C$加群として分裂する。
任意の$C$加群$N$は$B\otimes_CN$の
$C$直和因子となり、$B$-射影加群を制限すれば
$C$-射影的なので逆の不等式も得る。
\end{proof}

\begin{corollary}[可換中心ファイバーからの十分条件と一つの非正則値]
\label{cor:quotient-azumaya-fibers}
命題\ref{prop:quotient-azumaya}の仮定の下で、
$D_\gamma=R_1/(g-\gamma)$と置く。
$\gamma^2\notin\{0,-1,-28/27\}$なら
$\operatorname{gldim}D_\gamma=2$。
一方、$\gamma^2=-28/27$では
$\operatorname{gldim}D_\gamma=\infty$。
$\gamma=0$での正則性は別途
命題\ref{prop:local-model}の非可換$A_1$局所模型を用いる。
\end{corollary}
\begin{proof}
可換中心ファイバーの必要十分判別式
$H(H+1)(27H+28)$は
\href{https://github.com/uedakazushi/nc-cubic2/blob/main/papers/central_fiber_discriminant.tex}
{別の現行TeX}に証明がある。
$H=\gamma^2$が指定三値以外なら中心曲面は滑らか。
中心の正則元で割った商が正則局所環なら
元の中心局所環も正則なので、
各点で$R_1$はAzumayaであり、
$D_\gamma$は滑らかな二次元中心上Azumaya。
補題\ref{lem:quotient-azumaya-gldim}から結論する。

$\gamma^2=-28/27$では中心曲面の
$(U,V,W)=(1/189,0,0)$は通常二重点だが、
全中心の方程式の$\partial_g$は
$2\gamma\,\partial_HF_H\ne0$なので全中心は滑らか。
従ってこの点で$D_\gamma$は
特異な可換局所曲面上Azumayaとなり、
補題\ref{lem:quotient-azumaya-gldim}から無限大域次元。
\end{proof}


\section{中心の四つの通常二重点}
$k=\CC$とする。$A=R_1$を旧中心三次商のfiltered三変数代数、
$g=3G_1/2$を中心正則元とし、$C=Z(A)$と書く。
次の全中心表示を、命題\ref{prop:quotient-full-center}（有限語空間の有理数計算を実行確認済み、斉次中心定理は既知の入力）に基づいて用いる：
\begin{equation}\label{eq:center}
 C\simeq
 \CC[g,U,V,W]/(\Phi),\quad
 \Phi=49U^3-(13g^2+14)U^2+(g^2+1)^2U
 +(11U-g^2-1)VW+V^3-W^3 .
\end{equation}
また$A$が$C$上有限なCohen--Macaulay整域、中心的斜体の次数6、
Azumaya locusが$\Spec C$のsmooth locusに一致することは、上の中心の有限計算と既知SklyaninのAzumaya結果を入力にした命題\ref{prop:quotient-azumaya}で証明した。
これらは本稿の多項式検算のみから従う主張ではない。

\begin{lemma}\label{lem:four}
$\Spec C$の$g=0$上にある特異点は
\[
 p_0=(g,U,V,W)=(0,1/7,0,0),\qquad
 p_\zeta=(0,-1,-4\zeta,4\zeta^2)\quad(\zeta^3=1)
\]
の四点であり、各点の形式完備局所環は
\begin{equation}\label{eq:conifold}
 \widehat C_{p}\simeq
 R:=\frac{\CC[[p,q,r,\gamma]]}
 {pq-(r-\gamma)(r+\gamma)},\qquad (g)=(\gamma)
\end{equation}
という中心切断を保つ形式同型を持つ。
$C/(g)$の四点はすべて$A_1$型曲面特異点である。
\end{lemma}
\begin{proof}
$\Phi$およびその偏微分を$g=0$上で解くことにより
表示した四点が得られる。これが全解であることは
\eqref{eq:center}の$g=0$におけるGr\"obner消去、
または元の中心ファイバー判別式の計算で確認できる。
点$p_0$での$(g,U,V,W)$のHessian行列式は$-64/343$、
$p_1$では$576$であり、両者は非零である。
$\zeta$に関する$(V,W)\mapsto(\zeta V,\zeta^2W)$という
対称性により$p_\zeta$でも非零となる。
さらに$(U,V,W)$方向のHessianの行列式は
それぞれ$-32/7$, $288$である。
従って$g$を基底パラメータに固定した形式的陰関数定理と
パラメータ付きMorse補題により、
中心方程式は非退化三変数二次形式と
非零係数を持つ$g^2$の和に変換できる。
$g$の単元倍を$\gamma$と取り直し、三変数の二次形式を
$pq-r^2$に変換すると\eqref{eq:conifold}を得る。
$\gamma=0$は$pq-r^2=0$なので$A_1$である。
\end{proof}
\begin{remark}
中心の全体の正規性には、\eqref{eq:center}が整域であることと
特異locusが余次元2以上という$R_1$条件を用いる。
超曲面の$S_2$条件だけで正規性は従わない。
\end{remark}

% Source-recovered from the 2026-11-01 central-quotient manuscript, Sections 17--18.
% Input by papers/central_quotient_local_models.tex.
\section{中心値零の小次元表現と六次分裂元}
\label{sec:central-quotient-representations}

前節の代数$A=R_1$を
\[
 A=\CC\langle X,Y,Z\rangle/
 \bigl([X,Y]-Z^2,\ [Y,Z]-(X^2+1),\ [Z,X]-Y^2\bigr)
\]
とする。$\Omega_1=(X^3+Y^3+Z^3)/3
-(XYZ+XZY+YXZ+YZX+ZXY+ZYX)/9$、
$G_1=\Omega_1+X$、$g=3G_1/2$と置く。
本節の命題は旧論文VIの第17--18節から復元した。
関係式・$G_1$の明示表現上の消滅は、
\texttt{computations/verify\_central\_quotient\_reps.py}
により$\CC$の標数零の係数で厳密に検査できる。

\begin{proposition}[一次・二次元単純表現の不存在]
\label{prop:local-no-small-simples}
$g=0$上に一次元単純表現は存在せず、
$A$は二次元単純複素表現を持たない。
\end{proposition}
\begin{proof}
一次元表現では交換子は零なので
$Z=Y=0$、$X^2=-1$となる。
この場合$G_1=X^3/3+X=2X/3$で、
$g=X=\pm i\ne0$である。

二次元表現においてPauli行列$\boldsymbol{\sigma}$を使い
\[
 X=aI+\boldsymbol x\cdot\boldsymbol{\sigma},\quad
 Y=bI+\boldsymbol y\cdot\boldsymbol{\sigma},\quad
 Z=cI+\boldsymbol z\cdot\boldsymbol{\sigma}
\]
と書く。$\boldsymbol{x},\boldsymbol{y},\boldsymbol{z}\in\CC^3$とし
内積・外積は標準双線形形式を使う。
三つの関係式は
\[
 \begin{gathered}
 \boldsymbol x^2=-a^2-1,\quad
 \boldsymbol y^2=-b^2,\quad
 \boldsymbol z^2=-c^2,\\
 \boldsymbol x\times\boldsymbol y=-ic\boldsymbol z,\quad
 \boldsymbol y\times\boldsymbol z=-ia\boldsymbol x,\quad
 \boldsymbol z\times\boldsymbol x=-ib\boldsymbol y
 \end{gathered}
\]
と同値である。

$abc\ne0$なら三つのベクトルは互いに直交し、
$A_0=a^2,B_0=b^2,C_0=c^2$について
\[
 B_0(A_0+1)=C_0^2,\quad
 B_0C_0=A_0(A_0+1),\quad
 C_0(A_0+1)=B_0^2.
\]
第一・第三式から$(A_0+1)^2=B_0C_0$で、
第二式との比較は$A_0+1=A_0$を強制して矛盾する。
$abc=0$の場合は非零である$a,b,c$の個数で分ける。
二つ非零なら外積零と内積零により非零平方長のベクトルと
従属な零直交ベクトルから矛盾する。
$b$または$c$だけが非零の場合も
$\boldsymbol x^2=-1$と外積・内積の式から矛盾する。
$a$だけが非零の場合、$\boldsymbol x\ne0$なら
$\boldsymbol y,\boldsymbol z$はこれに比例し
$\boldsymbol y\times\boldsymbol z=-ia\boldsymbol x$に矛盾する。
$\boldsymbol x=0$なら$X=\pm iI$で、
$Y,Z$は互いに可換な平方零行列となり共通不変直線を持つ。
最後に$a=b=c=0$なら$X^2=-I$で
$[X,Y]=[Z,X]=0$だから$Y,Z$は$X$の固有基底で対角かつ
平方零、すなわち零である。この場合も可約。
従って二次元単純表現は存在しない。
\end{proof}

\subsection{軸上の二つの三次元単純表現}
\begin{proposition}[六次多項式による明示表現]
\label{prop:local-two-axial}
\[
 q(t)=7t^6+21t^4+14t^2+1,\quad
 \lambda_0=\frac{42t^5+119t^3+75t}{13},\quad
 \lambda_1=\frac{7t^5+35t^3+45t}{13},\quad
 \lambda_2=t
\]
とし、添字を$\mathbb Z/3$で読んで
$d_i=\lambda_{i+1}-\lambda_i$とする。
$q(t)=0$である任意の複素数$t$に対し
\[
 Xe_i=\lambda_i e_i,\qquad
 Ye_i=d_{i-1}e_{i+1},\qquad
 Ze_i=d_ie_{i-1}\qquad(i\in\mathbb Z/3)
\]
は$A/(G_1)$の三次元単純表現を与える。
$t$の六根から得られる同型類はちょうど二つであり、
$\operatorname{Tr}X=\pm\sqrt{-7}$で区別される。
両方の中心指標は$(g,U,V,W)=(0,1/7,0,0)$である。
両方において$Z^3$はスカラーであり、
二つの$Z$の固有値を合わせると$T^6+1/7$の六根となる。
\end{proposition}
\begin{proof}
上の行列に代入すれば三つの交換子関係と
$\Omega_1+X=0$が$\CC[t]/(q)$で成立する。
付属の厳密記号計算はすべての九行列成分の剰余零を検証する。
また
\[
 s=\lambda_0+\lambda_1+\lambda_2,\qquad
 s^2=-7,\quad
 d_0d_1d_2=-s/7,\quad
 Z^3=(-s/7)I
\]
が$q$での多項式剰余として成立する。
$q$は重根を持たず、$q$と$d_0d_1d_2$は互いに素である。
従って各$d_i\ne0$で、$X$は相異なる三固有値を持ち、
$Y$の三周期による推移性から表現は単純である。
$X$の三固有値は、$s$を固定したとき
$T^3-sT^2-2T+s/7$の三根である。
循環順序の違いは基底の循環置換で同一視される。
$s$は$\pm\sqrt{-7}$の二通りで、traceが異なる表現は同型でない。
二つでの$Z^3$はそれぞれ$\mp\sqrt{-7}/7$なので
固有値の六つの組は$T^6+1/7$の全ての根になる。
\end{proof}

\subsection{非軸上の三つの中心点を支配する行列}
原始三乗根$\omega^2+\omega+1=0$を固定し、
\[
 X_0=
 \begin{pmatrix}
 -1&1&1-\omega\\
 -1&(2\omega+1)/3&1\\
 -(\omega+2)/3&1/3&-(2\omega+1)/3
 \end{pmatrix},
 \quad Y_0=\operatorname{diag}(-1,\omega,\omega^2),
\]
\[
 Z_0=
 \begin{pmatrix}
 -1&-(2\omega+1)/3&\omega+1\\
 -\omega-1&(\omega-1)/3&\omega+1\\
 (\omega-1)/3&-\omega/3&-(\omega+2)/3
 \end{pmatrix}
\]
と置く。
\begin{proposition}[六つの非軸上表現]
\label{prop:local-six-nonaxial}
$\varepsilon=\pm1$、$k\in\mathbb Z/3$について
\[
 (X,Y,Z)=(\varepsilon X_0,\varepsilon\omega^kY_0,
                   \varepsilon\omega^{2k}Z_0)
\]
は$A/(G_1)$の三次元単純表現を与える。
その中心指標は旧論文VIの中心生成元$U,V,W$で
\[
 (g,U,V,W)=(0,-1,-4\omega^{2k},4\omega^k)
\]
であり、同じ中心点で二つの非同型類がある。
\end{proposition}
\begin{proof}
$\omega^2+\omega+1=0$上で
$[X_0,Y_0]=Z_0^2$、
$[Y_0,Z_0]=X_0^2+I$、
$[Z_0,X_0]=Y_0^2$および
$\Omega_1(X_0,Y_0,Z_0)+X_0=0$を直接確認できる。
これは付属の$\CC[\omega]/(\omega^2+\omega+1)$上の
有理数記号計算で再現する。
$Y_0$は相異なる三固有値を持ち、
$X_0$の非対角成分はすべて非零なので
共通不変部分空間は存在しない。
$g$は符号変換で反転し、$U,V,W$の規格化された
中心多項式の代入により表示した中心指標を得る。
$V,W$の非可換語による明示式を
\texttt{computations/verify\_central\_quotient\_reps.py}
に追加した。これらの中心指標は、
$\mathbb Q[\omega]/(\omega^2+\omega+1)$上の
独立した記号行列計算でも検証した。
なお保存版の全スクリプトの一括実行は未実施である。
異なる$k$は中心指標で、同じ$k$の符号は
$\operatorname{Tr}X=-\varepsilon$で区別できる。
\end{proof}

\begin{lemma}[六次最小多項式の回復]
\label{lem:local-z-minpoly}
$C=Z(R_1)$は系\ref{cor:quotient-normal-center}により正規整域であり、
中心元の定義は$U=Z^3(g-Z^3)$である。
このとき$K=\operatorname{Frac}C$上の
$Z$の最小多項式および被約特性多項式は
\[
 P(T)=T^6-gT^3+U
\]
である。
\end{lemma}
\begin{proof}
$P(Z)=0$なので最小多項式$m(T)$は$P(T)$を割る。
$m$の係数は$P$の根の対称式として$C$上整であり
$K$に属するため、$C$の正規性から$m(T)\in C[T]$。
軸上の二つの三次元表現は$g=0,U=1/7$という
同一の中心指標を持つ。ここで$U=-Z^6=1/7$は
軸上の$Z^3=-s/7$と$s^2=-7$から直接得られる。
二表現の$Z$固有値を併せると六つが相異なるので、
$m$をこの中心指標に特殊化すると六根を持つ。
従って$\deg m\ge6$であり、$m=P$。
次数六の中心的斜体での被約特性多項式は
次数六のモニック多項式で$m$の倍数だから同じものとなる。
\end{proof}

\begin{proposition}[八つの三次元単純表現の一次拡大]
\label{prop:quotient-full-ext1}
各四つの中心点の二つの三次元単純表現を$S_+,S_-$とする。
命題\ref{prop:local-model}のPBW・既知Sklyanin入力の下で
\[
 \dim\operatorname{Ext}^1_{R_1}(S_\epsilon,S_\eta)
 =\dim\operatorname{Ext}^1_{R_1/(g)}(S_\epsilon,S_\eta)
 =\begin{cases}0&\epsilon=\eta,\\2&\epsilon\ne\eta.\end{cases}
\]
異なる中心点に属する単純表現の間では、全てのExt群が零である。
従って一次拡大箙は、二頂点と各方向二矢印を持つ四つの非連結成分からなる。
\end{proposition}
\begin{proof}
命題\ref{prop:local-model}の完備化でMorita同値を取り、
$E=\operatorname{End}_R(R\oplus I)$、
$R=\CC[[p,q,a,b]]/(pq-ab)$、$I=(p,a)$とする。
ここで$\gamma=(b-a)/2$である。
$R$の極大イデアルを$\mathfrak m$とすると
\[
 I^*=p^{-1}(p,b),\qquad II^*=(p,a,b,q)=\mathfrak m.
\]
$R$と$I$は非同型直既約で、両者の自己準同型環は$R$。
従ってJacobson radicalは
\[
 J_E=\begin{pmatrix}\mathfrak m&I^*\\I&\mathfrak m\end{pmatrix},\qquad
 J_E/J_E^2\simeq
 \begin{pmatrix}0&I^*/\mathfrak m I^*\\I/\mathfrak m I&0\end{pmatrix}.
\]
$I$と$I^*$はそれぞれ二つの極小生成元を持つので、
二つの非対角成分はともに$\CC^2$である。
projective coverのradicalを使うExt$^1$の計算から、
自己拡大は零、相互拡大は二次元となる。

さらに$\gamma E\subset J_E^2$なので、$\gamma$で割っても
$J_E/J_E^2$は変化しない。$(g)=(\gamma)$であるから
中心切断$R_1/(g)$でも同じ一次拡大表を得る。
中心上有限な代数の有限長加群間のExtは同じ中心閉点に支持され、
忠実平坦完備化と可換なので、元の代数でのExtも同じである。
異なる中心指標を持つ二加群では、ある中心元が異なるスカラーで作用する。
Ext上の二つの中心作用は等しいため、その差の非零スカラーがExtを殺し、
従って全Ext群が零となる。

独立の有限計算でも同じ結果を確認した。
上三角block表現の非対角三行列を27変数として関係式を線型化した行列を$J$、
$g=Y^3-ZYX+X$の線型化を追加した行列を$J_g$とすると、
軸上・非軸上の両方について
\[
\begin{array}{c|c|c|c|c}
 &\operatorname{rank}J&\operatorname{rank}J_g
 &\operatorname{rank}(\text{内的微分})&\dim\operatorname{Ext}^1\\\hline
 \text{自己}&19&19&8&0\\
 \text{相互}&16&16&9&2
\end{array}
\]
となる。これは$\mathbb Q[t]/(q)$と
$\mathbb Q[\omega]/(\omega^2+\omega+1)$上の厳密な行列階数である。
再現コード\texttt{computations/verify\_central\_quotient\_extensions.py}は
$J$と$J_g$が内的微分を殺すことも検査する。
符号反転は各点の二表現を交換し、三乗根の対称性は非軸上三点を巡回するので、
この計算は八つの全単純表現を扱っている。
\end{proof}

\begin{proposition}[軸上の三次元表現の無限小剛性と整数小行列式]
\label{prop:quotient-axial-rigidity}
命題\ref{prop:local-two-axial}の二つの三次元単純表現$V$は、
$R_1$-加群としても$R_1/(G_1)$-加群としても
\[
 \operatorname{Ext}^1(V,V)=0
\]
を満たす。これは有限体上の非零$19$次小行列式により
特性零で証明できる。
\end{proposition}
\begin{proof}
$R_1$の三つの関係式を三行列$X,Y,Z$に代入した
$27$個の多項式を、各行列の$9$成分からなる$27$変数で
微分し、$27\times27$行列$J$を得る。
列は$X,Y,Z$それぞれに行優先で標準行列単位$E_{ij}$を並べ、
行は関係式
$[X,Y]-Z^2,\ [Y,Z]-X^2-I,\ [Z,X]-Y^2$
をそれぞれ行優先で並べる。

$q(t)=7t^6+21t^4+14t^2+1$は$\mathbb F_3$上既約であり、
従って$\mathbb Q$上既約である。
$\mathbb F_{29}$において$t=2$は$q$の根で、
\[
 (\lambda_0,\lambda_1,\lambda_2)=(3,10,2),\qquad
 (d_0,d_1,d_2)=(7,21,1)\pmod{29}.
\]
$J$の次の零始まり行番号・列番号からなる小行列を取る：
\[
\begin{split}
 I={}&(2,0,1,5,3,4,10,11,9,14,6,7,13,20,16,17,18,19,8),\\
 K={}&(0,1,2,3,4,5,6,7,8,9,10,11,12,13,14,16,18,19,20).
\end{split}
\]
整数係数の明示的なJacobi微分から計算すると
\[
 \det J_{I,K}\equiv15\not\equiv0\pmod{29}.
\]
再現コードは
\texttt{computations/verify\_central\_quotient\_reps.py}に
あり、各微分列を生成して法$29$のGaussian消去を行う。
従って$\mathbb Q[t]/(q)$上でrankは少なくとも$19$。
既約性により、すべての複素根の表現で同じ下界が成立する。

各$V$は単純なので共役軌道の次元は
$\dim\operatorname{PGL}_3=8$である。
軌道方向はJacobianの核に含まれるからrankは高々$19$。
従って接空間の次元は$27-19=8$で、軌道方向と一致する。
この接空間を内的導分で割ったものが
$\operatorname{Ext}^1_{R_1}(V,V)$なので零となる。
$G_1=0$を追加した表現スキームの接空間は前者の部分空間で
同じ共役軌道を含むため、こちらも次元は$8$であり、
$\operatorname{Ext}^1_{R_1/(G_1)}(V,V)=0$となる。
\end{proof}


\section{完備局所代数の分裂とconifoldのmodifying generator}
\begin{lemma}[既存の最大次数分裂元]\label{lem:splitting}
中心的斜体$A\otimes_C\Frac C$内の$Z$の最小多項式が
\begin{equation}\label{eq:degree-six}
 P(T)=T^6-gT^3+U
\end{equation}
であり、被約特性多項式も同式であることは補題\ref{lem:local-z-minpoly}によって導かれる。
四点$p$において$\Lambda=A\otimes_C\widehat C_p$は
rank-one corner $e\Lambda e=\widehat C_p$を持ち、
\[
 \Lambda\simeq\End_R(M),\quad
 R\mid M,\quad M\in\CM(R),\quad \operatorname{rank}_R M=6.
\]
\end{lemma}
\begin{proof}
四点では$g=0$かつ$U\ne0$なので、$P(T)$の剰余は
$T^6+U(p)$であり、六つの根は相異なる。
Henselの補題によって$\widehat C_p$上の六つの根
$\lambda_i$に持ち上がり、全ての差が単元となる。
従って
\[
 e_i=\prod_{j\ne i}\frac{Z-\lambda_j}{\lambda_i-\lambda_j},
 \qquad 1=\sum_{i=1}^6e_i
\]
が直交冪等元の分解を与える。
被約特性多項式の次数が6なので、各$e_i$の一般行列rankは1。
$e=e_1$とし$M=\Lambda e$と置けば、
$\Lambda$のMCM性により$M$と$e\Lambda e$は反射的である。
高さ1ではAzumaya性から$e\Lambda e=R$、また
$\Lambda\to\End_R(M)$は同型であり、
反射性によってそれぞれ全体へ延長される。
$e\Lambda e=R$が$M$の直和因子であることと
rankの等式も従う。
\end{proof}
\begin{remark}[分裂元に必要な実証]
旧論文VIでは$U=Z^3(g-Z^3)$と二つの三次元単純表現の
相異なる$Z$固有値を用いて、\eqref{eq:degree-six}の
最小次数が6であることを証明している。
軸上の二つの明示表現とその全行列は第\ref{sec:central-quotient-representations}節にTeX化し、
関係式を\texttt{computations/verify\_central\_quotient\_reps.py}で
厳密に独立検算した。中心全生成元と$U$の恒等式は第\ref{sec:quotient-center-azumaya}節、
中心の正規性は系\ref{cor:quotient-normal-center}で確立した。
\end{remark}
\begin{proposition}[非可換$A_1$局所模型]\label{prop:local-model}
補題\ref{lem:splitting}の仮定に加え、
命題\ref{prop:local-no-small-simples}の小次元単純表現不存在を用いる。
このとき、\eqref{eq:conifold}の座標を選んで
\begin{equation}\label{eq:endR}
 \widehat A_p\simeq\Mat_3\bigl(\End_R(R\oplus I)\bigr),
 \qquad I=(p,r-\gamma)
\end{equation}
が成り立つ。$I=(p,r+\gamma)$としても同型の記述が得られる。
\end{proposition}
\begin{proof}
$R$はconifold型の三次元完備局所超曲面である。
$M$は$R$を直和因子に持ち、$\End_R(M)$はMCMなので
modifying generatorである。
Iyama--WemyssのcAn型局所分類
\cite[Theorem~5.1(1)]{IW}を適用する。
同定理の基底は$\CC[[r,\gamma]]$、
$f=(r-\gamma)(r+\gamma)$の二因子は相異なる素元である。
$M$はMCMで$R$を直和因子に持ち、$\End_R(M)$もMCMなので、
同定理のmodifying generatorの仮定を満たす。
因子集合$\{1,2\}$のflagは空flag、$\{1\}$、$\{2\}$だけなので、
基本generatorは$R$または
$R\oplus(p,r\pm\gamma)$に限る。
よって$M\simeq R^a\oplus I^b$、$a\ge1$と書ける。
$b=0$なら$\End_R(M)$はAzumayaとなり、
命題\ref{prop:quotient-azumaya}の中心特異点で非Azumayaという結論に矛盾するので$b\ge1$。
二つの非同型直和因子のend環の半単純剰余は
$\Mat_a(\CC)\oplus\Mat_b(\CC)$である。
低次元単純表現の不存在から$a,b\ge3$。
一方$\operatorname{rank}_R M=6$より$a+b=6$なので
$a=b=3$となり、\eqref{eq:endR}を得る。
\end{proof}

\section{中心切断とBrauer指数}
\begin{proposition}[局所切断のskew group表示]\label{prop:skew}
$S=R/(\gamma)=\CC[[p,q,r]]/(pq-r^2)$、
$J=(p,r)\subset S$と置くと
\[
 \End_R(R\oplus I)/(\gamma)
 \simeq\End_S(S\oplus J)
 \simeq\CC[[a,b]]\rtimes\{\pm1\},
\]
ただし非自明な元は$(a,b)\mapsto(-a,-b)$と作用する。
従って\eqref{eq:endR}の中心切断は
\[
 \widehat{(A/(g))}_{p}\simeq
 \Mat_3\bigl(\CC[[a,b]]\rtimes\{\pm1\}\bigr)
\]
であり、大域次元は2である。
\end{proposition}
\begin{proof}
$I^\vee=\Hom_R(I,R)=p^{-1}(p,r+\gamma)$であり、
$\End_R(R\oplus I)$は
$\begin{psmallmatrix}R&I^\vee\\ I&R\end{psmallmatrix}$である。
$R/I\simeq\CC[[q,\gamma]]$なので
$I\cap\gamma R=\gamma I$、従って$I/\gamma I\simeq J$。
$I^\vee$についても同様であり、行列積を保ったまま
$\End_S(S\oplus J)$を得る。
$T=\CC[[a,b]]$の符号反転作用に対し
$S=T^{\{\pm1\}}$であり、
$(p,q,r)=(a^2,b^2,ab)$と置ける。
偶奇分解$T=S\oplus T^{-}$において
$T^{-}\simeq J$であるため
$\End_S(T)\simeq T\rtimes\{\pm1\}$となる。
有限群の位数は可逆なので、このskew group algebraの
大域次元は$T$と同じ2である。
\end{proof}

\begin{corollary}[零ファイバーの局所正則性と単純表現]\label{cor:zero}
命題\ref{prop:quotient-azumaya}のPBW・既知Sklyanin入力と
上記の局所模型を用いる（小単純表現排除は命題\ref{prop:local-no-small-simples}で証明した）。
$D_0=A/(g)$の中心は$C/(g)$であり、
$\gldim D_0=2$であり、$D_0$はhomologically smoothな二次元Calabi--Yau代数である。
各四つの$A_1$中心点には三次元単純表現が二つ、
それ以外の中心閉点には六次元単純表現が一つ存在する。
従って$D_0$の三次元単純表現は全部で8同型類となる。
\end{corollary}
\begin{proof}
まず$C\cap gA=gC$である。実際、$ga\in C$なら
$g[a,x]=0$と$g$の正則性から$a\in C$となる。
従って$C/(g)\hookrightarrow A/(g)$であり、商はこの中心部分環上有限である。
中心の滑らかな点ではAzumaya性により
完備化は$\Mat_6(\CC[[a,b]])$となる。
四特異点では命題\ref{prop:skew}の模型がある。
両方の模型の中心は$C/(g)$の対応する完備局所環に等しい。
中心の形成を有限個の交換子写像の核として表し、
中心基底への忠実平坦完備化と有限生成性を使って
$Z(D_0)=C/(g)$を得る。
各閉点の完備局所代数の大域次元は2なので、
有限生成加群の第二syzygyの射影性を忠実平坦降下して
$\gldim D_0\le2$を得る。
滑らかな中心点の単純加群の第二Extは非零であるから等号成立。
命題\ref{prop:quotient-constant-cy}と
補題\ref{lem:tails-central-cy}により、有限大域次元と中心有限性から
homological smoothnessと二次元CY性も従う。
剰余半単純代数は特異点で
$\Mat_3(\CC)\oplus\Mat_3(\CC)$、
滑らかな中心点で$\Mat_6(\CC)$であり、
単純表現の分類が従う。
\end{proof}

\begin{corollary}[局所分裂と大域Brauer障害]\label{cor:brauer}
上記と同じPBW・既知Sklyanin入力を用いる。
このとき$D_0$は整域で、一般中心分数体$K_0$上で
$D_0\otimes_{Z(D_0)}K_0$は次数6の中心的斜体となり、
Brauer指数は6である。
特に$D_0$は大域的な$\End_{Z(D_0)}(N)$型に分裂しない。
一方、四つの中心特異点の完備化では
命題\ref{prop:skew}の行列模型へ分裂する。
\end{corollary}
\begin{proof}
PBW同型の下で$g$の最高次記号は非零正則元$3\Omega_1/2$であり、
$\deg(ag)=\deg a+3$から
$\operatorname{gr}(A/(g))=S_1/(\Omega_1)$となる。
右辺は滑らかな楕円曲線のtwisted homogeneous coordinate ringで、
その関数体のskew Laurent環に埋め込まれる
\cite[Lemmas~2.18, 2.19]{WWYSklyanin}から整域である。
従って$D_0$も整域となる。
中心有限な$D_0$の一般中心局所化は
次数6、次元36の有限次元中心的単純代数であり、
整域性から斜体となる。そのBrauer指数は次数と等しく6。
$\End(N)$なら一般点では行列環となって矛盾する。
局所完備化での分裂は一般中心分数体のBrauer類を
零にするものではない。
\end{proof}

\section{零中心ファイバーのCayley三次曲面と楕円境界}
\label{sec:quotient-cayley}
\begin{proposition}[選ばれた反標準因子を持つCayley表示]
$C_0=Z(D_0)$の射影閉包$\overline S\subset\mathbb P^3_{U,V,W,T}$は
\[
 P=49U^3-14U^2T+UT^2+(11U-T)VW+V^3-W^3=0
\]
で定まり、Cayley三次曲面である。
$\omega^2+\omega+1=0$として
\[
 \begin{pmatrix}U\\V\\W\\T\end{pmatrix}
 =\begin{pmatrix}
 -9&-1&-1&-1\\
 0&-4&-4\omega&-4\omega^2\\
 0&4&4\omega^2&4\omega\\
 -63&1&1&1
 \end{pmatrix}
 \begin{pmatrix}a_0\\a_1\\a_2\\a_3\end{pmatrix}
\]
と置くと、変換は可逆で
\[
 P=-1728\sum_{i<j<k}a_i a_j a_k,\qquad
 T=-63a_0+a_1+a_2+a_3.
\]
$\operatorname{Spec}C_0=\overline S\setminus E$で、
$E=\overline S\cap\{T=0\}$は滑らかな反標準楕円曲線である。
\end{proposition}
\begin{proof}
変換行列の行列式は$-3456(2\omega+1)\ne0$。
$\omega^2+\omega+1$で剰余を取る直接代入で二式を得る。
無限遠の曲線は
\[
 E:\quad49U^3+11UVW+V^3-W^3=0.
\]
その三偏微分が零なら、いずれかの座標が零のとき全座標が零となる。
全座標が非零なら$147U^2=-11VW$、$3V^2=-11UW$、$-3W^2=-11UV$を掛け、
$1323=1331$という矛盾を得る。
従って$E$は滑らかで、無限遠に曲面の特異点もない。
三次超曲面の随伴公式$\omega_{\overline S}=\mathcal O_{\overline S}(-1)$から反標準性が従う。
\end{proof}

\begin{corollary}[Brauer指数6のtwisted NCCR]
$D_0$は正規Gorenstein曲面$C_0$のtwisted noncommutative crepant resolutionである。
そのBrauer指数は6であり、通常の$\operatorname{End}_{C_0}(N)$型のNCCRではない。
\end{corollary}
\begin{proof}
$C_0$は四つの$A_1$点以外で滑らかで、超曲面なので正規Gorenstein。
$R_1$の中心Macaulay性を中心正則元$g$で割ると$D_0$は$C_0$上MCMである。
従って反射的で、高さ1ではAzumaya、かつ大域次元2である。
\cite[Definitions~2.1, 2.2]{VdBOverview}のtwisted NCCRの条件を満たす。
一般中心分数体上は次数6の斜体なので、通常の自己準同型環なら行列環となって矛盾する。
局所完備化での非捩れ模型と、この大域Brauer障害を区別する。
\end{proof}
\begin{remark}[非可換affine del Pezzoとの比較]
この中心曲面は、選んだ滑らかな反標準楕円曲線を除いたCayley三次曲面であり、
\cite[\S1.1, Theorems~3.4.4, 3.4.5]{EGDelPezzo}の非可換affine del Pezzoの具体的な有限PI例と比較できる。
同文献のgenericな一生成元中心を、この特殊PI値の全中心へ適用してはいけない。
また定数項変形のPBW・CY性自体は\cite[Theorem~3.3.2]{EGDelPezzo}の既知の一般定理の特殊化であり、
本稿では具体的なsyzygyと自己双対解消も与えた。
Brauer類の具体的なsymbol表示、大域stack上の実現と大域対象の復元は未解決である。
\end{remark}

\section{古典的pencil、量子的判別式と非被約な点加群}
\label{sec:quotient-classical-comparison}
\begin{proposition}[古典的pencilと三様体の判別式]
\[
 F_\alpha=\frac{x^3+y^3+z^3}{3}+\alpha uvx+u^3+v^3
\]
とし、$v=au$の曲面切断を考える。
有限pencilパラメータでの特異条件と三様体の特異条件はそれぞれ
\[
 \Delta_\alpha(a)=4\alpha^3a^3+9(1+a^3)^2=0,\qquad
 F_\alpha\text{が特異}\ \Longleftrightarrow\ \alpha^3=-9.
\]
$p(x)=x^3/3+\alpha ax+1+a^3$に対し
$\operatorname{Res}_x(p,p')=\Delta_\alpha(a)/9$で、無限遠切断は滑らかである。
さらに
\[
 \operatorname{Disc}_a\Delta_\alpha
 =2^{12}3^{14}\alpha^9(\alpha^3+9)^3,\qquad
 \Delta_\alpha(\omega)=4(\alpha^3+9)\quad(\omega^3=1).
\]
\end{proposition}
\begin{proof}
$u=0$上のFermat平面三次曲線は滑らかであり、$u=1$の特異点では$y=z=0$。
従って曲面では$p,p'$の共通根を調べればよい。
三様体の勾配は
$y=z=0$、$x^2+\alpha uv=0$、$\alpha vx+3u^2=0$、$\alpha ux+3v^2=0$を与える。
射影点では$u,v\ne0$で、$u=1$の消去は
$3x-\alpha^2v^2=v^3-1=\alpha^3+9=0$となる。
従って三特異点は$[\alpha^2\omega^2/3:0:0:1:\omega]$である。
残りは一変数resultantとdiscriminantの厳密計算から従う。
\end{proof}
\begin{remark}[重複度を含むresultantの規格化]
$\operatorname{Res}(X_1^2,\ldots,X_n^2)=1$で規格化すると
\[
 \operatorname{Res}(\nabla F_\alpha)=3^8(\alpha^3+9)^{12},\qquad
 \operatorname{Res}(\nabla(F_\alpha|_{v=au}))=\Delta_\alpha(a)^4.
\]
最初は$(x,u,v)$の三変数resultant $9(\alpha^3+9)^3$の四乗である。
$x=1$における$\CC[u,v]/(3u^2+\alpha v,3v^2+\alpha u)$の基底$1,u,v,uv$を使うと、
$1+\alpha uv$の乗法行列式は$(\alpha^3+9)^3/729$となる。
最高次係数$3^8$とPoisson積公式\cite[equation~(2)]{DKSResultant}で三変数の式を得る。
切断の式も、残る二変数の勾配resultantと$y^2,z^2$の二つの平方から従う。
これらは古典的判別式であり、量子的なtails判別式ではない。
\end{remark}

\begin{proposition}[正則性を区別しない定常点の粗い局所データ]
$h\ne0$、$\alpha\ne0$、$\alpha^3\ne-9$とする。
$A_{\alpha,h}$の中心一次元が非零となるpoint moduleのスキームは、
六点を台とし、各局所環が$\CC[y,z]/(y^2,z^2)$である長さ24の有限スキームとなる。
\end{proposition}
\begin{proof}
可換試験環上で$u$が可逆なら、中心性の多重線型関係
$u_i a_{i+1}=a_i u_{i+1}$から$u_i=1$に規格化した点列は一定となる。
従ってこのchartではpoint moduleは可換化で表される。
可換化の関係は$y^2=z^2=0$、$x^2+\alpha v=0$である。
$v=-x^2/\alpha$を消去すると中心三次式は
\[
 1-\frac{2x^3}{3}-\frac{x^6}{\alpha^3}-\frac{2h^2}{3}xyz=0.
\]
$y=z=0$での$x^3$の方程式は
$3q^2+2\alpha^3q-3\alpha^3=0$で、判別式は$4\alpha^3(\alpha^3+9)$。
二根は相異なり非零なので六つの$x$を与える。
$x$方向の微分は各点で単元であり、形式的陰関数定理で$x$を$yz$の式として消せる。
従って各局所環は表示した長さ4の環である。
$u=0,v\ne0$の可換化では$x^2=y^2=z^2=0$かつ$v^3=0$となり矛盾するので、
中心座標が非零の点は全てこのchartに含まれる。
\end{proof}
\begin{remark}
$h=1$で$\alpha^3=-243/28$なら、古典三様体と楕円境界は滑らかでもtailsは非正則である。
$\alpha=1$の正則例とこの三つの非正則例は、上の定常点の台の個数、
局所環型と全長が一致する。埋込みを含む全point schemeの同型を主張してはいない。
$\alpha\ne0,\alpha^3\ne-9$の古典pencilには六つの特異曲面があり、
$\alpha=0$では判別式は$9(a^3+1)^2$、三切断は滑らかなFermat曲線上の錐である。
前者の各特異曲面は、$x$方向の非退化二次項を消去すると
局所Jacobian環$\CC[y,z]/(y^2,z^2)$（長さ4）を持つ。
$\alpha=0$の三つの錐の頂点では$\CC[x,y,z]/(x^2,y^2,z^2)$となり長さ8である。
全切断の滑らかさやpencil判別式の重根の有無を、三様体自体の滑らかさと混同しない。
\end{remark}


\section{二中心変数を逆にしたchartの21個の中心特異点}
\label{sec:quantum-boundary-center}
$\alpha^3=-243/28$とし、旧論文VIの二つの中心変数$u,v$を逆にした
次数零chartを$Q=A_{\alpha,1}[u^{-1},v^{-1}]_0$、
$t=v/u\in\CC^\times$とする。
\begin{equation}\label{eq:Ht}
 \begin{aligned}
 H(t)&=-\frac7{27}(t^3+2+t^{-3}),\\
 F_H&=49U^3-(13H+14)U^2+(H+1)^2U\\
 &\quad +(11U-H-1)VW+V^3-W^3 .
 \end{aligned}
\end{equation}
\begin{proposition}[二重\'etale被覆による中心の忠実平坦降下]
\label{prop:chart-descent}
系\ref{cor:quotient-ambient-pbw}の相対PBW性・中心正則性を用いる。
中心表示\ref{prop:quotient-full-center}の下で
\begin{equation}\label{eq:chart-center}
 Z(Q)\simeq
 \CC[t,t^{-1},U,V,W]/(F_{H(t)}).
\end{equation}
また$Q$は中心上有限で一般rank36である。
比較写像の有限平坦性と分岐位置は以下の明示式から従い、
独立した未証明の入力とはしない。
\end{proposition}
\begin{proof}
$Q$に$s^2=\alpha t$を添加し、$x/u=sX,y/u=sY,z/u=sZ$と置く。
$s,\alpha,t$は可逆だから、代入は可逆であり、
\[
 Q'=Q\otimes_{\CC[t,t^{-1}]}\CC[s,s^{-1}]
 \simeq R_1\otimes_{\CC[g]}\CC[s,s^{-1}],\qquad
 g\longmapsto-\tfrac32\left(s^{-3}+s^3/\alpha^3\right).
\]
左の拡大は階数2の有限忠実平坦\'etale拡大である。
右辺の係数拡大は
\[
 \CC[s,s^{-1}]
 \simeq\CC[g,s]/\left(s^6+\tfrac23\alpha^3gs^3+\alpha^3\right)
\]
というmonic表示を持ち、定数項が可逆なので$s$も可逆。
従って$\CC[g]$上自由rank6である。
微分$\frac92(s^{-4}-s^2/\alpha^3)$は
$s^6\ne\alpha^3$、すなわち$t^3\ne1$で可逆となる。
さらに
\[
 g(s)^2-\frac9{\alpha^3}
 =\frac{9(s^6-\alpha^3)^2}{4\alpha^6s^6}
\]
だから、この開集合は基底の$g^2-9/\alpha^3$を逆にしたものに等しく、
比較はそこで有限\'etaleである。

有限個の代数生成元との交換子写像の核は中心なので、
平坦基底変更は中心の形成と可換である。
従って$Z(Q')=Z(R_1)\otimes_{\CC[g]}\CC[s,s^{-1}]$。
被覆の対合は$s\mapsto-s$、$(X,Y,Z)\mapsto(-X,-Y,-Z)$で、
$g$を反転し$U,V,W$を固定する。
$g(s)^2=H(t)$なので、この中心は
$\CC[s,s^{-1},U,V,W]/(F_{H(t)})$である。
$\CC[s,s^{-1}]$は$\CC[t,t^{-1}]$上の
基底$1,s$を持ち、対合は後者の成分を反転する。
不変部分を取れば\eqref{eq:chart-center}を得る。
有限性と一般rank36もこの忠実平坦被覆で検出できる。
なお比較構成とmonic表示は任意の$\alpha\ne0$で有効であり、
その場合は$H(t)=9(t^3+2+t^{-3})/(4\alpha^3)$と置けばよい。
\end{proof}

\begin{proposition}[中心の21個の特異閉点]\label{prop:21}
\eqref{eq:chart-center}の中心特異点は21個に尽き、次の通りである：
\[
\begin{array}{c|c|c|c}
t\text{の条件}&H(t)&(U,V,W)&\text{点数}\\ \hline
t^3=-1&0&(1/7,0,0),\,(-1,-4\zeta,4\zeta^2),\ \zeta^3=1&12\\
7t^6-13t^3+7=0&-1&(0,0,0)&6\\
t^3=1&-28/27&(1/189,0,0)&3
\end{array}
\]
$H=0$と$H=-28/27$に対応する15点は三次元通常二重点である。
\end{proposition}
\begin{proof}
$F_H$の三つの偏微分を同時に零とする
中心ファイバーは$H=0,-1,-28/27$に限る。
各ファイバーの特異点は表の通りである。
これらは$F_H$と$\partial_UF_H,\partial_VF_H,\partial_WF_H$の
Gr\"obner基底で検算できる。全中心の$t$微分は
$(\partial_HF_H)H'(t)$であるが
\[
 H'(t)=-\frac79(t^2-t^{-4}).
\]
$H=0$と$H=-28/27$では$t^6=1$、
$H=-1$の点$(U,V,W)=(0,0,0)$では
$\partial_HF_H=0$なので、表の全点が全中心でも特異である。
それ以外に候補はなく、$t$の各方程式は相互に共通根を持たない。
$7s^2-13s+7$の判別式は$-27$であり、$s=0$の根を持たないので
それぞれの根は三つの$t$を与え、点数は$12+6+3=21$である。
通常二重点性も厳密に検算できる。
局所座標$q=t^3$を使うと$H(q)=-7(q+2+q^{-1})/27$で、
$q=\pm1$では$H'=0$、$H''=-14/(27q^3)\ne0$。
従って四変数Hessianの行列式は
\[
 (\partial_HF_H)H''(q)\,
 \det\operatorname{Hess}_{U,V,W}(F_H).
\]
$q=-1$では、軸上点と非軸上点でそれぞれ
$\partial_HF_H=1/49,1$、三変数行列式$-32/7,288$である。
$q=1$では$\partial_HF_H=-1/1323$、
三変数行列式$-8/1701$であり、いずれも非零。
$dt\mapsto dq=3t^2dt$は可逆なので、15点すべての通常二重点性が従う。
残る6点では四変数Hessianはrank1である。
これらは\texttt{computations/verify\_central\_quotient\_local\_hessian.py}で
有理数および$\CC[q]/(7q^2-13q+7)$上の厳密計算として再現できる。
\end{proof}

\begin{theorem}[ホモロジー的特異点の同定：明示的平坦比較]
\label{thm:three-homological}
命題\ref{prop:quotient-azumaya}のPBW・既知Sklyanin入力と
命題\ref{prop:chart-descent}の明示的比較の下で、
$Q$の有限大域次元を持たない中心閉点は
\[
 \mathfrak p_\tau=(t-\tau,U-1/189,V,W),\qquad \tau^3=1
\]
の三点に限る。各点の完備局所代数は
\[
 \widehat Q_{\mathfrak p_\tau}\simeq
 \Mat_6\!\left(
   \CC[[a,b,c,d]]/(a^2+b^2+c^2+d^2)
 \right).
\]
従って中心の残りの18特異点は、
中心環が非正則でも代数はホモロジー的に正則である。
\end{theorem}
\begin{proof}
命題\ref{prop:chart-descent}の比較は
$s^2=\alpha t$の被覆での
$\CC[g]\to\CC[s,s^{-1}]$という写像に基づく。
この写像の分岐値は$g^2=-28/27$、逆像は$t^3=1$である。
従って$t^3\ne1$では有限\'etale中心基底変換によって
大域次元3の代数から$Q$を得る。
有限\'etale基底変換で大域次元上界が保たれるのは、
制限・誘導とseparability idempotentを用いればよい。
逆向きは有限忠実平坦降下により
有限生成加群の第3 syzygyの射影性を検出する。
ゆえに$t^3\ne1$での局所代数の大域次元は高々3。

$t^3=1$に対応する全中心上では、
元の$R_1$は中心上Azumayaである。
実際、中心曲面の唯一の特異点では$\partial_g\Phi\ne0$で、
他の点では$(U,V,W)$方向が既に滑らかなので、
全中心のこのファイバーは全点が滑らかである。
ここで命題\ref{prop:quotient-azumaya}を適用する。
忠実平坦降下から$Q$もAzumayaである。
$\mathfrak p_\tau$以外の中心点は滑らかなので
局所大域次元は有限である。
三点の中心は三次元通常二重点である。
完備局所中心上のAzumaya代数を$L$とすると、
剰余環$L/\mathfrak mL=\Mat_6(\CC)$のprimitive idempotentを
完備性により$e\in L$へ持ち上げられる。
$eLe$は中心上rank1の自由加群で、単位写像は剰余体上同型だから中心と同型。
$e$はfullで$Le$は自由rank6なので、
$L\simeq\operatorname{End}(Le)\simeq\Mat_6(\widehat Z)$となる。
形式Morse補題で$\widehat Z$を表示したnodeに直せばよい。通常二重点の剰余体は無限射影次元なので、
この三点の局所代数は有限大域次元を持たない。
\end{proof}
\begin{remark}[特異圏の同値の証明の所在]
三つのホモロジー的特異点からKaroubi完備な射影特異圏
$\operatorname{Perf}_{\mathbb Z/2}(\CC)^{\oplus3}$
への条件付き同値は、後の
第\ref{sec:quotient-singcat}節で二つの中心局所化と
Iwanaga--Gorenstein完備化によって証明する。
中心有限性とAzumaya locusは第\ref{sec:quotient-center-azumaya}節、
有限平坦比較は命題\ref{prop:chart-descent}に証明した。
\end{remark}

% 現行研究資料：旧中心商特異性論文VI 第21--22節からの回収
% 分類結果の固有の証明と外部入力（とりわけ有限中心・Azumaya・正則性）を区別する。
\section{三つの特異点の射影特異圏}
\label{sec:quotient-singcat}

本節では$A_\alpha$を中心三次商の一混合項族
\[
 A_\alpha=\frac{\CC[u,v]\langle x,y,z\rangle}
 {([x,y]-z^2,\,[y,z]-(x^2+\alpha uv),\,[z,x]-y^2,\,
   \Omega_1+\alpha uvx+u^3+v^3)}
\]
とする。ただし
\[
 \Omega_1=\frac{x^3+y^3+z^3}{3}
       -\frac{xyz+xzy+yxz+yzx+zxy+zyx}{9}
\]
であり、$u,v$は一次中心元である。
相対PBW性と周囲のAS正則性は
系\ref{cor:quotient-ambient-pbw}で証明した。
中心有限性・Azumaya locusと二重被覆の比較は
命題\ref{prop:quotient-azumaya}、\ref{prop:chart-descent}で証明した。
以下で用いる境界の正則性とIwanaga--Gorenstein性も、
次の補題でこれらの入力から導く。

\begin{lemma}[境界の正則性とchartのGorenstein性]
\label{lem:quotient-boundary-gorenstein}
$\alpha\ne0$とする（比較構成は命題\ref{prop:chart-descent}の末尾の一般化を用いる）。
$A_\alpha/(u)$のqgrは大域次元2を持ち、
$u$-chartの$t=v/u=0$での商$S_1/(\Omega_1+1)$は
大域次元高々2を持つ。
また$Q_\alpha=A_\alpha[u^{-1},v^{-1}]_0$は
左右noetherianで中心上有限なIwanaga--Gorenstein代数であり、
左右自己入射次元は高々3である。
\end{lemma}
\begin{proof}
$A_\alpha/(u)=S_1[v]/(\Omega_1+v^3)$であり、
$v$-chartは$S_1/(\Omega_1+1)$、境界は$S_1/(\Omega_1)$である。
Euler導分$\delta(\Omega_1)=3\Omega_1$と
補題\ref{lem:gluing-moving-central-value}でchartの大域次元は高々2。
境界は滑らかな楕円曲線のcoherent圏なので、
命題\ref{prop:gluing-one-chart}と
補題\ref{lem:gluing-serre-dimension}からqgrの大域次元は2となる。

命題\ref{prop:chart-descent}の階数2の中心\'etale被覆で
\[
 Q'_\alpha\simeq
 R_1[s,s^{-1}]/\left(g+\tfrac32(s^{-3}+s^3/\alpha^3)\right).
\]
$R_1[s,s^{-1}]$は左右noetherianで大域次元高々4。
括弧内は非零中心正則元なので、中心正則元に対するReesの補題は
任意の有限生成左右$Q'_\alpha$加群$M$について
\[
 \operatorname{Ext}^i_{Q'_\alpha}(M,Q'_\alpha)
 \simeq
 \operatorname{Ext}^{i+1}_{R_1[s,s^{-1}]}
                  (M,R_1[s,s^{-1}])
\]
を与える。従って左右自己入射次元は高々3。
中心平坦基底変更で有限生成加群のExtは可換であり、
忠実平坦性がその消滅を検出するので、この上界は$Q_\alpha$へ降下する。
中心上有限性・noetherian性は命題\ref{prop:chart-descent}による。
\end{proof}

\newcommand{\SingCat}{\operatorname{D}_{\mathrm{sg}}}
\newcommand{\PerfCat}{\operatorname{Perf}}
\newcommand{\thickcat}{\operatorname{thick}}
\newcommand{\Kar}{\mathrm{Kar}}
\newcommand{\Db}{\mathrm D^b}
$\mathcal P_\alpha=\thickcat\{\pi A_\alpha(n):n\in\mathbb Z\}
\subset\Db(\operatorname{qgr}A_\alpha)$と置き、
\[
 \mathcal S_\alpha=
 \Db(\operatorname{qgr}A_\alpha)/\mathcal P_\alpha
\]
を射影的特異圏とする。上付き$\Kar$はKaroubi完備化を表す。

\begin{lemma}[中心元に沿う特異圏の局所化]
\label{lem:quotient-central-localization}
$H$は左右noetherianな次数付き代数、
$z\in Z(H)_1$は正則元とする。
$\Db(\operatorname{qgr}(H/(z)))$が構造対象の全twistで
厚く生成されるなら、中心局所化は
\[
 \left(\Db(\operatorname{qgr}H)/
       \thickcat\{\pi H(n)\}\right)^\Kar
   \simeq \SingCat(H[z^{-1}]_0)^\Kar
\]
を与える。次数なしの左右noetherian代数$\Lambda$と
中心正則元$z$に対し、$\Lambda/(z)$の大域次元が有限なら
同じ結論が成立する。
\end{lemma}
\begin{proof}
中心局所化で零となる有限生成加群は
$z$のある冪で殺される。$z$の冪による有限filtrationの
各商は$H/(z)$上のcoherent加群なので、仮定と
$0\to H(-1)\xrightarrow{z}H\to H/(z)\to0$
によって局所化の導来核は$\thickcat\{\pi H(n)\}$に入る。
局所化後の有限生成加群には有限個の生成元の
$H$-spanというlatticeが存在する。
有界複体の微分と射・homotopyの有限個の分母も
$z$の冪で払えるので、有界導来圏の局所化は
$z$-torsion cohomologyを零にするVerdier局所化となる。
twistsは局所化後に自由加群となり、逆に有限自由複体は
twistsの有限和とconeから得られる。
局所化後のperfect対象はこれらの直和因子なので、
直和因子の持上げを仮定せず両辺をKaroubi完備化すればよい。
非次数付きの場合は有限次元torsionによる商を省く。
\end{proof}

\begin{corollary}[二つの中心chartの除去]
\label{cor:quotient-two-chart}
補題\ref{lem:quotient-boundary-gorenstein}の下で
\[
 \mathcal S_\alpha^\Kar\simeq\SingCat(Q_\alpha)^\Kar .
\]
\end{corollary}
\begin{proof}
$A_\alpha/(u)\simeq
 S_1[v]/(\Omega_1+v^3)$は
$v$の非零スカラー変換で既知の正則曲面のモデルに同型である。
同補題の大域次元二とtwistsの厚い生成性により、
まず補題\ref{lem:quotient-central-localization}を$z=u$に適用する。
$u$-chartにおける$t=0$の商は
$S_1/(\Omega_1+1)$であり、
Euler導分$\delta(\Omega_1)=3\Omega_1$を用いる
補題\ref{lem:gluing-moving-central-value}により有限大域次元を持つ。
同補題のaffine版を$t$に適用して$Q_\alpha$に至る。
共通境界$u=v=0$も最初のfiltrationに含まれる。
\end{proof}

\begin{lemma}[孤立行列nodeの完備化と直和分解]
\label{lem:quotient-completion}
$C$をaffine可換noetherian環、
$\Lambda$を$C$上有限な左右noetherian Iwanaga--Gorenstein代数とする。
有限個の極大イデアル$\mathfrak m_i$を除いて
$\Lambda_{\mathfrak p}$が有限大域次元であり、
各完備化が$\widehat\Lambda_{\mathfrak m_i}\simeq
\operatorname{Mat}_{d_i}(N_i)$で、
$N_i$は孤立通常二重点の完備局所環とする。
さらに各局所nodeの単純剰余加群を与える
$\Lambda/\mathfrak m_i\Lambda$の有限加群が存在するとする。
このとき
\[
 \SingCat(\Lambda)^\Kar\simeq
 \bigoplus_i\SingCat(N_i)^\Kar .
\]
\end{lemma}
\begin{proof}
Iwanaga--Gorenstein性に対するBuchweitz--Happel--Chenの同値
\cite{ChenRelativeSingcat}により、
左辺は有限生成Gorenstein-projective加群のstable圏で計算できる。
$M,N$がそのような加群ならstable Homは
\[
 \underline{\operatorname{Hom}}_\Lambda(M,N)=
 \operatorname{coker}\!\left(
 N\otimes_\Lambda\operatorname{Hom}_\Lambda(M,\Lambda)
 \longrightarrow\operatorname{Hom}_\Lambda(M,N)\right)
\]
で、有限$C$加群である。
例外点以外では$M$はGorenstein-projectiveかつ
有限射影次元なので射影的である。
従ってstable Homは有限個の閉点だけに支持されて有限長。
有限表示性、中心の平坦完備化、有限長加群の局所分解により、
各Hom群は各完備化のstable Homの直和に同型となる。
同じ議論はshiftにも適用でき、完備化は完全忠実である。

本質的全射性をKaroubi完備化後に証明する。
通常二重点$N_i=k[[z_1,\ldots,z_e]]/(q)$の
maximal Cohen--Macaulay加群を行列因子化で表すと、
$d^2=q$の微分から$\partial_{z_j}q$は特異圏で零作用となる。
非退化二次式ならこれは$z_j$を零化するので、
$z_j$のKoszul複体は元の対象を直和因子として含み、
そのcohomologyは剰余体の有限直和である。
従ってnodeの特異圏は剰余体で厚く生成される。
仮定した有限$\Lambda$加群はその生成元を完備化の像として与え、
直和因子まで閉じれば所望の圏全体になる。
\end{proof}

\begin{theorem}[三つの通常二重点に集中する射影特異圏]
\label{thm:quotient-three-singcat}
$28\alpha^3+243=0$とし、
上記のPBW・既知Sklyanin入力を用いる。
中心局所化の前提とIwanaga--Gorenstein性は
補題\ref{lem:quotient-boundary-gorenstein}、
三つの局所模型とそれ以外の点での有限大域次元は
定理\ref{thm:three-homological}から従う。
このとき
\[
 \mathcal S_\alpha^\Kar\simeq
 \bigoplus_{\tau^3=1}
  \SingCat\!\left(
   \frac{\CC[[a,b,c,d]]}{a^2+b^2+c^2+d^2}
  \right)^\Kar
 \simeq \PerfCat_{\mathbb Z/2}(\CC)^{\oplus3}.
\]
従って$K_0(\mathcal S_\alpha^\Kar)\simeq\mathbb Z^3$で、
非零直既約対象は同型を除いて六つである。
$[1]$は各二つの直既約対象を交換し、$[2]\simeq\mathrm{id}$。
\end{theorem}
\begin{proof}
系\ref{cor:quotient-two-chart}で$Q_\alpha$に局所化する。
局所模型\ref{thm:three-homological}と
補題\ref{lem:quotient-completion}から最初の圏同値を得る。
ここでは中心上有限な代数の剰余環が$\Mat_6(\CC)$なので、
そのcolumn単純加群を$Q_\alpha$加群と見れば、
補題\ref{lem:quotient-completion}の生成元の存在条件も満たされる。

二次元node $N_0=k[[p,q]]/(pq)$のMCM加群は
正規化$k[[p]]\oplus k[[q]]$上の飽和を取り、
conductorと剰余体の線形部分空間を標準化することにより、
直和因子として$N_0$, $k[[p]]$, $k[[q]]$からなる。
stable圏では$N_0$は零、後の二つは相互直交で
stable Endが$\CC$、shiftは二つを交換する。
よって$\SingCat(N_0)^\Kar\simeq\PerfCat_{\mathbb Z/2}(\CC)$。
まずaffineな$k[p,q]/(pq)$と$k[p,q,r,s]/(pq+rs)$に
OrlovのKn\"orrer周期性\cite[Theorem~2.1]{OrlovSingcat}を適用する。
両方の特異点は原点だけなので、補題\ref{lem:quotient-completion}を
可換環に適用して完備化へ移せる。従って
$N_0$と$k[[p,q,r,s]]/(pq+rs)$の
Karoubi完備な特異圏は同値である。
$\CC$上で後者の二次式は四変数の平方和と線形同型なので結論。
二周期的な一因子の$K_0$は$\mathbb Z$で、
三因子の直和が残りの主張を与える。
\end{proof}
\begin{remark}
全ての同値はKaroubi完備化\emph{後}の主張である。
$D^b(\operatorname{qgr}A_\alpha)$そのものや
非可換cubicのKuznetsov成分と同値であるとは主張しない。
中心の21個の特異点と、代数のホモロジー的特異点三個を区別する。
\end{remark}

\section{正則三次商の球面ツイストと構造対象の標識}
\label{sec:quotient-spherical-twist}
本節は$\alpha=1$の正則例を対象とする。
$A=A_1$、$\mathcal O=\pi A$とし、
正則性$\operatorname{gldim}\operatorname{qgr}A=3$は
定理\ref{thm:quotient-exact-tails}から従う。
命題\ref{prop:quotient-ambient-nakayama}と
補題\ref{lem:gluing-serre-dimension}の双対性から
$\mathbb S_{\Db(\operatorname{qgr}A)}\simeq(-2)[3]$も従う。
さらにfiltered三変数代数
\[
 R_1=
 \frac{\CC\langle X,Y,Z\rangle}
 {([X,Y]-Z^2,\ [Y,Z]-(X^2+1),\ [Z,X]-Y^2)},
 \quad G_1=\Omega_1(X,Y,Z)+X
\]
に、命題\ref{prop:local-two-axial}から得られる三次元単純左加群の線形双対である\emph{右}加群$V$を取る。この加群は
\[
 \dim_\CC V=3,\qquad VG_1=0,\qquad
 \operatorname{End}_{R_1}(V)=\CC,\quad
 \operatorname{Ext}^1_{R_1}(V,V)=0
\tag{QS}
\]
を満たす。単純性と$VG_1=0$は命題\ref{prop:local-two-axial}、無限小剛性は命題\ref{prop:quotient-axial-rigidity}で証明される。後者の法$29$の厳密小行列式は現在の検算コードから再現できる。左加群から右加群への双対は有限次元自己拡大と単純性を保つ。

$t=v/u$とし、
$Q^\circ=A[u^{-1},v^{-1},(v^3-u^3)^{-1}]_0$と置く。
$t=-1$、$(\bar x,\bar y,\bar z)=(iX,iY,iZ)$
（$i^2=-1$）により$V$を$Q^\circ$-加群とする。
次数付き右$A$-加群$N_V$を
\[
 (N_V)_n=
 \begin{cases}V,&n\ge0,\\0,&n<0,\end{cases}
 \qquad u=\mathrm{id},\ v=-\mathrm{id},\
 x=iX,\ y=iY,\ z=iZ
\]
という次数一作用で定め、$\mathcal T=\pi N_V$と書く。
$\mathcal T$は非零coherent対象で、
$u$がtailで可逆なので$\mathcal T(1)\simeq\mathcal T$。

\begin{lemma}[局所自己拡大の消滅]
\label{lem:quotient-rigid-chart}
(QS)を用いると
\[
 \operatorname{End}_{Q^\circ}(V)=\CC,\qquad
 \operatorname{Ext}^1_{Q^\circ}(V,V)=0 .
\]
さらに命題\ref{prop:quotient-azumaya}と
命題\ref{prop:chart-descent}の$\alpha=1$での比較から
$\operatorname{gldim}Q^\circ\le3$。
\end{lemma}
\begin{proof}
自己拡大$0\to V\to E\to V\to0$に対し、
中心作用$N=t+1$は$N^2=0$である。
$s=i+N/(2i)$と置くと$s^2=t$。
$X=\bar x/s$, $Y=\bar y/s$, $Z=\bar z/s$と置けば
$E$は$R_1$-加群になり、
\[
 G_1=-b(s),\qquad b(s)=s^{-3}+s^3
\]
を満たす。
(QS)により$R_1$-自己拡大は分裂するため
$G_1$は$E$上でも零。
$b(i)=0$、$b'(i)=-6$、$(s-i)^2=0$から
$s=i$、$t=-1$を得る。
分裂は$Q^\circ$-線形となり、Extが零である。
大域次元については$s^2=t$を添加すると
\[
 g=-\frac32(s^3+s^{-3}),\quad
 g^2-9=\frac{9(s^6-1)^2}{4s^6}
\]
となる。$t^3-1$を逆にした領域では
$s\mapsto g$の微分は単元で、
$R_1[(g^2-9)^{-1}]$からの有限\'etale中心基底変換となる。
有限\'etale拡大での大域次元保存と忠実平坦降下から
上界三を得る。
\end{proof}

\begin{lemma}[coherent対象への導来随伴]
\label{lem:quotient-chart-ext}
上記の仮定の下で
\[
 \mathbf R\operatorname{Hom}(\mathcal O,\mathcal T)
   \simeq V[0],\quad
 \operatorname{End}(\mathcal T)=\CC,\quad
 \operatorname{Ext}^1(\mathcal T,\mathcal T)=0 .
\]
また$\mathcal T$は大きい$\operatorname{QGr}A$内で
長さ高々三のinjective分解を持つ。
\end{lemma}
\begin{proof}
$A^\circ=Q^\circ[u,u^{-1}]$と書ける。
中心局所化から
$j^*:\operatorname{QGr}A\to\operatorname{Mod}Q^\circ$
を得る。その右随伴
$j_*M=\pi(M\otimes_{Q^\circ}A^\circ)$は、
$A^\circ$が左$Q^\circ$加群として自由なのでexactである。
ここで局所化済み加群$L=M\otimes_{Q^\circ}A^\circ$はtorsion-freeかつ飽和している。
実際、$A_{\ge n}\to L$という次数付き写像は、中心元$u^n$の像を
$u^{-n}$倍することで$A\to L$へ一意に延長する。
$V$から作る非負次数格子$N_V$に対して$L/N_V$はtorsionである
（各負次数元は十分高い正次数で非負次数へ入る）。
従って$j_*V\simeq\pi N_V=\mathcal T$である。
$j^*$もexactだから導来随伴
\[
 \mathbf R\operatorname{Hom}_{\operatorname{QGr}A}
 (F,j_*V)\simeq
 \mathbf R\operatorname{Hom}_{Q^\circ}(j^*F,V)
\]
が成立する。$j^*\mathcal O=Q^\circ$、
$j^*\mathcal T=V$を代入して前半を得る。
$\operatorname{QGr}A$は局所noetherianなGrothendieck圏で、
coherent対象間のExtは$\operatorname{qgr}A$内と一致する。
最後に$\operatorname{gldim}Q^\circ\le3$で
$V$の長さ三以下のinjective分解を取り、
exactな$j_*$で移せばよい。
\end{proof}

\begin{theorem}[構造対象をline objectのshiftに送らない自己同値]
\label{thm:quotient-spherical-counterexample}
上記で証明した$\alpha=1$の正則性とSerre関手を用いる。
(QS)は命題\ref{prop:local-two-axial}、\ref{prop:quotient-axial-rigidity}により成立する。$\mathcal T$は3-spherical対象であり、
Seidel--Thomasの球面ツイスト
\cite[Definition~2.14, Proposition~2.10]{SeidelThomas}
\[
 \mathbb T_{\mathcal T}(F)=
 \operatorname{Cone}\left(
 \mathbf R\operatorname{Hom}(\mathcal T,F)
    \otimes_\CC\mathcal T\longrightarrow F
 \right)
\]
は$\Db(\operatorname{qgr}A)$の自己同値を与える。
標準$t$構造に関して
\[
 \mathcal H^j(\mathbb T_{\mathcal T}(\mathcal O))=
 \begin{cases}
 \mathcal O,&j=0,\\
 \mathcal T^{\oplus3},&j=2,\\
 0,&j\ne0,2.
 \end{cases}
\]
従って全ての$m,r\in\mathbb Z$について
$\mathbb T_{\mathcal T}(\mathcal O)\not\simeq\mathcal O(m)[r]$。
\end{theorem}
\begin{proof}
$\mathcal T(1)\simeq\mathcal T$だから
$\mathbb S\mathcal T\simeq\mathcal T[3]$。
補題\ref{lem:quotient-chart-ext}とSerre双対から
\[
 \operatorname{Ext}^j(\mathcal T,\mathcal T)=
 \begin{cases}\CC,&j=0,3,\\0,&j\ne0,3.\end{cases}
\]
負次数はheartの二対象間のExtで消え、
高次数は$\operatorname{gldim}=3$により消える。
Serre双対性は必要な合成対の非退化性を与える。
補題の有限injective分解と
coherent対象間の有限次元Extにより
Seidel--Thomasの球面対象の条件を満たす。
その際、大きい圏$\operatorname{QGr}A$は局所noetherianな
Grothendieck圏なので十分なinjectiveを持ち、injectiveの直和もinjectiveである。
小さい圏$\operatorname{qgr}A$はnoetherian対象からなるSerre部分圏で、
全ての大きい対象はcoherent部分対象の有向和である。
従ってcoherent対象への全射にはcoherent部分対象からの全射が存在し、
Seidel--Thomasの条件(C1)--(C4)を満たす。
有限大域次元とproper性により、各coherent対象との全Extも有限次元である。
これでinjective enhancement上の評価写像とconeの構成が適用できる。
これは同論文の一般Abel圏での構成であり、
可換スキーム上の対象に限定しない。

さらにSerre双対と上記随伴から
\[
 \mathbf R\operatorname{Hom}(\mathcal T,\mathcal O)
  \simeq
  \mathbf R\operatorname{Hom}
     (\mathcal O,\mathbb S\mathcal T)^*
  \simeq V^*[-3].
\]
よってツイストの三角形は
$\mathcal T^{\oplus3}[-3]\to\mathcal O
\to\mathbb T_{\mathcal T}(\mathcal O)
\to\mathcal T^{\oplus3}[-2]$となり、
cohomology長完全列から表示が従う。
$\mathcal O$と$\mathcal T$はどちらも非零なので、
二つの異なる次数に非零cohomologyを持ち、
単一coherent対象のshiftである$\mathcal O(m)[r]$とは同型でない。
\end{proof}
\begin{remark}[標識なし復元とは異なる]
\label{rem:quotient-unmarked-reconstruction}
この自己同値は「任意の導来同値が構造対象のtwist/shiftを
保つ」という自動性を否定する。
異なる二つの代数の標識なし導来圏が同値なら
次数付き同型かどうか、また球面ツイスト等を合成すれば
標識を戻せるかどうかまでは決定しない。
従って構造対象を指定した復元定理とは矛盾しない。
\end{remark}


% 中心chartと導分によるperfect性の局所大域判定。
\section{中心chartと導分による射影的perfect性の判定}
\label{sec:quotient-chart-gluing}
以下の判定は有限位数・無限位数の楕円境界を区別しない。
使う外部入力は、noetherian AS正則代数の中心正則超曲面に対する
行列因子化と特異圏の同値（右加群には反対環を適用）
\cite[Theorems~4.7, 5.5]{CCKM}である。
中心元のnormalizing automorphismは恒等なので、
同文献のtwisted因子化は次数shiftを伴う通常の二周期因子化となる。

\begin{lemma}[有限射影因子化の縮約と分母払い]
\label{lem:gluing-contraction}
$2$が可逆な左右noetherian環$S$と中心正則元$w$を取り、
有限生成射影加群の行列因子化を$d^2=w$と書く。
$S/(w)$上のcokernelが有限射影次元なら、この因子化は縮約可能である。
さらに中心正則元$u$で局所化した後に縮約可能なら、
ある$N$について$u^N\mathrm{id}=dK+Kd$となる。
次数付きの場合も、因子化の次数shiftを含めて同じ主張が成立する。
\end{lemma}
\begin{proof}
$\bar S=S/(w)$上の二周期分解から、cokernelの有限射影次元は
各周期syzygyの射影性を意味する。従って分解は分裂し、
$\bar S$上で$d\bar H+\bar Hd=1$となる奇写像を取れる。
有限射影性からこれを$S$上の$H$へ持ち上げると
$dH+Hd=1+wT$となる。
左辺は$d$と可換なので、$w$の正則性により$[d,T]=0$。
ゆえに$H'=H-\frac12dT$は$dH'+H'd=1$を満たす。
局所化後の$H'$の有限個の成分の分母を$u^N$で払い、
$S\hookrightarrow S[u^{-1}]$の単射性を使えば最後の等式を得る。
\end{proof}

$B$を標数零の体上の左右noetherian connected graded AS正則代数、
$W\in Z(B)$を正次数の正則斉次元、$A=B/(W)$とする。
\[
 \mathcal P_A=\operatorname{thick}\{\pi A(n):n\in\mathbb Z\},\qquad
 \mathcal Q_A=\mathrm D^b(\operatorname{qgr}A)/\mathcal P_A
\]
と置く。以下で消滅させるのは$\mathcal Q_A$であり、
次数付き代数そのものの$\mathrm D_{\mathrm{sg}}^{\mathrm{gr}}(A)$ではない。

\begin{proposition}[一つの中心chartと正則境界]
\label{prop:gluing-one-chart}
$u\in Z(B)_1$が$B$と$A$上正則であり、
\[
 \mathrm D^b(\operatorname{qgr}A/(u))
 =\operatorname{thick}\{\pi(A/(u))(n):n\in\mathbb Z\}
\]
とする。$A[u^{-1}]_0$が有限大域次元を持つなら、
$\mathcal Q_A=0$である。
\end{proposition}
\begin{proof}
coherent対象のgraded liftを$M$とする。
$u:M(-1)\to M$のconeのcohomologyは$u$に殺される。
境界の生成性と
$0\to A(-1)\xrightarrow{u}A\to A/(u)\to0$
により、このconeは$\mathcal P_A$に属する。
従って$u$の作用は$\mathcal Q_A$で可逆である。
一方、\cite[Theorems~4.7, 5.5]{CCKM}から$M$の
次数付き特異圏の像を有限graded自由行列因子化で表せる。
$A[u^{-1}]\simeq A[u^{-1}]_0[u,u^{-1}]$のgraded加群圏は
次数零部分の加群圏と同値なので、局所化後のcokernelは有限射影次元を持つ。
補題\ref{lem:gluing-contraction}で縮約の分母を払うと、
$u^N$の作用は次数付き特異圏で零となり、従って$\mathcal Q_A$でも零となる。
可逆かつ零なので$M$の像は零である。
全coherent対象に適用して結論する。指数$N$は対象ごとに選べばよい。
\end{proof}

\begin{proposition}[中心因子を持つ導分による切断判定]
\label{prop:gluing-derivation}
次数零の真の導分$\delta:B\to B$が存在し、
\[
 \delta(W)=c\prod_{j=1}^r l_j,\qquad c\ne0,
\]
で、各$l_j\in Z(B)$は正次数の斉次元で$A$上正則とする。
各$A/(l_j)$のqgrの有界導来圏がその構造対象のtwistsで厚く生成されるなら、
$\mathcal Q_A=0$である。
\end{proposition}
\begin{proof}
前命題のconeの議論により、各$l_j$は$\mathcal Q_A$で可逆に作用する。
一方、有限graded自由行列因子化$d^2=W$の行列成分へ$\delta$を作用させると
\[
 d\,\delta(d)+\delta(d)\,d=\delta(W)\mathrm{id}
\]
なので、$\delta(W)$の作用は次数付き特異圏で零である。
従って積$\prod_j l_j$は$\mathcal Q_A$で零かつ可逆となり、全対象が零となる。
\end{proof}

\begin{lemma}[中心値を動かす導分]
\label{lem:gluing-moving-central-value}
$R$は左右noetherianで$\operatorname{gldim}R\le n<\infty$（$n\ge1$）、
$w=g+\beta$は中心正則元とする。
導分$\delta$が$\delta(g)=P(g)$、$P(-\beta)\ne0$を満たせば
$\operatorname{gldim}R/(g+\beta)\le n-1$である。
\end{lemma}
\begin{proof}
有限生成$D=R/(w)$加群$M$の$(n-1)$次syzygy $N$は
$R$上の射影次元が高々1である。
実際、$D$-自由分解の各短完全列と$\operatorname{pd}_R D=1$から
次元を順に下げればよい。
有限射影分解$0\to P_1\xrightarrow{\phi}P_0\to N\to0$を取り、
$wN=0$から$\phi\psi=w$、単射性から$\psi\phi=w$を得る。
各$P_i$は有限自由加群の直和因子なので$\delta$-connectionを持つ。
その和を$\nabla$とすると$D(d)=[\nabla,d]$は$R$-線形で
$dD(d)+D(d)d=P(g)$である。
$c=P(-\beta)$、$P(g)-c=wQ(g)$と置けば
\[
 H=c^{-1}\left(D(d)-\tfrac12Q(g)d\right),\qquad dH+Hd=1.
\]
従って$N$は$D$上射影的で$\operatorname{pd}_D M\le n-1$。
noetherian環の巡回加群による大域次元判定から結論する。
\end{proof}

\begin{lemma}[perfect性とSerre双対性からの正確な次元]
\label{lem:gluing-serre-dimension}
$A$は左右noetherian connected graded AS--Gorensteinで、
balanced dualizing complexが$A_\nu(-\ell)[d+1]$（$d\ge0$）で、
$\pi A\ne0$とする。
$\mathrm D^b(\operatorname{qgr}A)=\mathcal P_A$なら
$\operatorname{gldim}\operatorname{qgr}A=d$である。
\end{lemma}
\begin{proof}
\cite[Theorem~4.2, Corollary~4.3]{YZSerre}は左右の$\chi$と
有限cohomological dimensionを与える。
特に$\pi A(n)$は有限射影次元を持つ。
Chen\cite[Theorem~3.10]{ChenGeneralizedSerre}の一般化Serre双対性は
有限射影次元の対象に適用でき、その関手は
$S(M)=M_\nu(-\ell)[d]$となる。
仮定により全coherent対象が適用範囲に入るので、
\[
 \operatorname{Ext}^i(M,N)^*
 \simeq\operatorname{Ext}^{d-i}(N,M_\nu(-\ell)).
\]
$i>d$では右辺が負次数ゆえ零。
また$M=\pi A$、$N=M_\nu(-\ell)$と取ると
$\operatorname{Ext}^d(M,N)^*\simeq\operatorname{Hom}(N,N)\ne0$なので等号となる。
\end{proof}


% 分離型判別式、超越的中心値、有限位数境界の開な正則範囲。
\section{一般の中心零化元と二つのchart}
\label{sec:quotient-global-regularity}
本節は指定した一混合項族だけでなく、中心三次商の一般的な十分条件を扱う。
二次ambientの一般構成そのものを分類するものではなく、以下に明記した
AS正則性・良いPBW chart・滑らかな量子的境界を満たす固定周囲に適用する。
$B$を左右noetherian connected graded AS正則代数、
$W\in Z(B)$を正次数の正則斉次元、$A=B/(W)$とする。
$\mathcal P_A=\operatorname{thick}\{\pi A(n)\}$、
$\mathcal Q_A=\mathrm D^b(\operatorname{qgr}A)/\mathcal P_A$と書く。

\begin{lemma}[有限大域次元からtwistsによる生成へ]
\label{lem:global-finite-to-perfect}
$A$は次数1生成で、$\operatorname{qgr}A$が有限大域次元を持つなら、
$\mathrm D^b(\operatorname{qgr}A)=\mathcal P_A$である。
\end{lemma}
\begin{proof}
大域次元を$e$とし、graded自由表示を繰り返して
$0\to K\to L_e\to\cdots\to L_0\to M\to0$を取る。
各$L_i$は$\pi A(n)$の有限直和である。
最後のYoneda拡大は$\operatorname{Ext}^{e+1}(M,K)=0$なので、
$M$は有限複体$[L_e\to\cdots\to L_0]$の直和因子である。
line object自体が射影的であるとは仮定していない。
\end{proof}

\begin{proposition}[中心零化元による十分条件]
\label{prop:global-annihilators}
$u,v\in Z(B)_1$が$A$上正則列をなし、
$\mathrm D^b(\operatorname{qgr}A/(u,v))$が構造対象のtwistsで厚く生成されるとする。
正次数の斉次中心元$p_1,\ldots,p_r\in Z(A)$が
次数付き特異圏の全対象を零化し、ある$N$について
$(u,v)^N\subset(p_1,\ldots,p_r)A$なら、$\mathcal Q_A=0$である。
\end{proposition}
\begin{proof}
graded lift $M$のKoszul複体$K(p_1,\ldots,p_r;M)$のcohomologyは
$(p_1,\ldots,p_r)$に殺され、従って$(u,v)$の冪に殺される。
有限filtrationの各商は$A/(u,v)$加群である。
正則列のKoszul分解から$\pi(A/(u,v))(n)\in\mathcal P_A$なので、
境界の生成性によりこのKoszul複体は$\mathcal P_A$に入る。
一方、$\mathrm D^b(\operatorname{gr}A)\to\mathcal Q_A$は
次数付きperfect複体を殺すから次数付き特異圏を経由する。
従って各$p_i$は$\mathcal Q_A$で零作用し、反復coneは
$M$のshiftとtwistの直和へ分裂する。ゆえに$M$自身も零となる。
\end{proof}

\begin{proposition}[二つの中心chartの貼合せ]
\label{prop:global-two-charts}
$u,v\in Z(B)_1$が$B$上正則で$A$上正則列をなし、
$A[u^{-1}]_0,A[v^{-1}]_0$が有限大域次元を持つとする。
$\operatorname{qgr}A/(u,v)$の有界導来圏が構造対象のtwistsで厚く生成されるなら、
$\mathcal Q_A=0$である。
\end{proposition}
\begin{proof}
各coherent対象を行列因子化で表し、各chart上で縮約する。
補題\ref{lem:gluing-contraction}の分母払いから、対象ごとに選んだ
$u^r,v^s$が$\mathcal Q_A$で零作用する。
二重Koszul coneのcohomologyは$(u^r,v^s)$に殺され、
$(u,v)^{r+s-1}\subset(u^r,v^s)$から境界の有限filtrationを持つ。
前命題と同じ議論でconeはperfectである。
中心作用は自然変換なので、最初の零coneの直和分解とも可換し、
二重coneは元の対象の四つのshift・twistの直和となる。
従って元の対象は零である。一様な$r,s$は要求しない。
\end{proof}

\section{分離型二変数三次式の必要十分判別式}
\label{sec:quotient-separated-discriminant}
\begin{theorem}[分離族の正則性と二変数三次判別式]
\label{thm:global-separated}
$S$を次数1生成、Hilbert級数$(1-t)^{-3}$の左右noetherian
三次元AS正則代数、$g\in Z(S)_3$を正則元とする。
$\operatorname{qgr}(S/(g))\simeq\operatorname{coh}E$で、$E$は滑らかな楕円曲線、
構造対象のtwistsはampleな直線束列と仮定する。
$u,v$を中心一次変数として
\[
 A_f=S[u,v]/(g+f),\qquad f=au^3+bu^2v+cuv^2+dv^3
\]
と置く。このとき
\[
 \operatorname{gldim}\operatorname{qgr}A_f
 =\begin{cases}3,&\Delta_f\ne0,\\\infty,&\Delta_f=0,\end{cases}
\quad
 \Delta_f=b^2c^2-4ac^3-4b^3d-27a^2d^2+18abcd.
\]
\end{theorem}
\begin{proof}
$B=S[u,v]$はAS正則次元5、index 5であり、$W=g+f$は中心正則元である。
$S$次数のfiltrationから
$\operatorname{gr}_S A_f=(S/(g))[u,v]$を得るので、$u,v$は中心正則列。
$A_f$はAS--Gorenstein次元4、index 2である。

$\Delta_f\ne0$なら、$f_u,f_v$の共通零点は原点のみで
$(f_u,f_v)$は$(u,v)$-primaryである。
次数2としての斉次resultantは$\operatorname{Res}(f_u,f_v)=-3\Delta_f$。
$B$の真の導分$\partial_u,\partial_v$は
$d\,\delta(d)+\delta(d)d=\delta(W)$から$f_u,f_v$を特異圏の零化元にする。
楕円境界の生成性と命題\ref{prop:global-annihilators}でperfect性を得て、
補題\ref{lem:gluing-serre-dimension}により正確な次元3となる。

$\Delta_f=0$なら$\mathrm{GL}_2$変換とスカラーの吸収により
$f=u^2v,u^3,0$のいずれかにできる。
$v$-chartは順に$S[w]/(g+w^2),S[w]/(g+w^3),(S/(g))[w]$であり、
全てaugmentationを持つ。
最初の環には$\deg S_1=2,\deg w=3$、後二つには通常次数を入れる。
Hilbert級数の逆数は順に
\[
 \frac{(1-t^2)^3}{1+t^3},\qquad
 \frac{(1-t)^4}{1-t^3},\qquad
 \frac{(1-t)^4}{1-t^3}.
\]
これらはそれぞれ原始6乗根、原始3乗根で極を持ち、多項式ではない。
従って自明加群$k$のminimal graded自由分解は無限である。
これをungraded分解として$\operatorname{Hom}(-,k)$を適用しても
全微分が零なので、ungradedの射影次元も無限となる。
一方$M=A_f/(S_+,u)\simeq\CC[v]$のtailはchart上で$k$へ移る。
有限大域次元を仮定すれば補題\ref{lem:global-finite-to-perfect}により
このtailはperfect、従ってchart上の$k$もperfectとなって矛盾する。
\end{proof}
\begin{corollary}[分離族の四係数と同型類]
$\Delta_f\ne0$を満たす全ての$A_f$は、$S$を固定して$u,v$の線形変換を許すと同型である。
\end{corollary}
\begin{proof}
$\mathbb P^1$の相異なる三点の組は$\mathrm{PGL}_2$で移り合い、
残るスカラーは二変数の共通スカラーの三乗に吸収できる。
従って分離型四係数は、座標を忘れた四次元モジュライを与えない。
\end{proof}

\section{超越的中心値と固定ambientの一般中心三次商}
\label{sec:quotient-generic-central-values}
\begin{lemma}[超越的中心値の有限大域次元]
\label{lem:global-transcendental}
$k_0\subset\CC$、$R_0$は$k_0$-代数、$z\in Z(R_0)$は$k_0$上超越的とする。
$R=R_0\otimes_{k_0}\CC$が左右noetherianで$\operatorname{gldim}R\le n$なら、
$k_0$上超越的な全ての$\lambda\in\CC$について
$\operatorname{gldim}R/(z-\lambda)\le n$である。
\end{lemma}
\begin{proof}
$\lambda$を含む$\CC/k_0$の超越基底$T$を選び、$K=k_0(T\setminus\{\lambda\})$とする。
体拡大は自由かつ忠実平坦で、$K$-線形な$\CC\to K$の$1\mapsto1$という分裂がある。
従って大域次元の上界とnoetherian性は$R_K=R_0\otimes K$へ降下する。
$z$の冪は$K$上でも独立なので、中心局所化
$R_\eta=R_K\otimes_{K[z]}K(z)$も大域次元高々$n$。
$z\mapsto\lambda$で$K(z)=K(\lambda)$とみなすと
$R_\eta\otimes_{K(z)}\CC=R/(z-\lambda)$である。
最後の体拡大は代数的分離拡大で、右辺は$R$の商ゆえnoetherian。
\cite[Lemmas~3.2, 3.3]{RSVBertini}の仮定が揃い、上界が保たれる。
任意の超越的体拡大が有限大域次元を保つと仮定したわけではない。
\end{proof}

\begin{theorem}[固定した良いambientでのvery-generalな正則性]
\label{thm:global-very-general}
$B$は中心一次元$u,v$を持つ左右noetherian AS正則五次元代数とし、
$G\in Z(B)_3$、$\operatorname{gr}_{x,y,z}B=S\otimes\CC[u,v]$とする。
$S$は整域、$\operatorname{gr}G=\gamma\in Z(S)_3$は非零正則元とする。
全ての$f\in\CC[u,v]_3$について$B/(G+f)$は$\CC[u,v]$上自由であり、
$\operatorname{qgr}(S/(\gamma))$は滑らかな楕円曲線のcoherent圏と仮定する。
有限表示の定義体$k_0$に対し、$a,d$が$k_0(b,c)$上代数的独立なら
\[
 \operatorname{gldim}\operatorname{qgr}
 B/(G+au^3+bu^2v+cuv^2+dv^3)=3.
\]
\end{theorem}
\begin{proof}
$D_u=B[u^{-1}]_0$はnoetherianで有限大域次元を持つ。
実際、中心局所化は大域次元を増やさず、
$B[u^{-1}]=D_u[u,u^{-1}]$から制限・誘導と評価$u=1$で上界を降下できる。
$t=v/u$とすると商の$u$-chartは
$D_u/(G/u^3+bt+ct^2+dt^3+a)$。
括弧内の$a$を除く中心元は$k_0(b,c,d)$上定義され、
最高主変数成分が非零正則なので、その冪は独立である。
前補題でchartの有限大域次元を得る。
$v$-chartも$d$を中心値として同様に扱える。
命題\ref{prop:global-two-charts}と
補題\ref{lem:gluing-serre-dimension}が結論を与える。
\end{proof}

\subsection{任意の量子パラメータでの真の混合例}
まず任意の$h\in\CC$について
\begin{gather*}
 S_h=\CC\langle X,Y,Z\rangle/
 ([X,Y]-hZ^2,[Y,Z]-hX^2,[Z,X]-hY^2),\\
 \Omega_h=\frac{X^3+Y^3+Z^3}{3}
 -\frac{h^2}{9}\sum_{\sigma\in S_3}X_{\sigma(1)}X_{\sigma(2)}X_{\sigma(3)}
\end{gather*}
と置く。$S_h$は全ての有限$h$でKoszul AS正則三次元である。
これはSklyaninパラメータ$[1:-1:-h]$に対する標準定理
（\cite{ATVRegular,ISSklyanin}、後者Theorem~1.1）による。
非正則例外条件$p^3=q^3=r^3$は$1\ne-1$により生じない。
滑らかな楕円境界を用いるのは以下で$h^6\ne-27$に制限してからである。
前節のPBW証明は$r_i=[X_{i+1},X_{i+2}]-hX_i^2$、$\beta(r_X)=h m$にそのまま適用できる。
従って$[y,z]=h(x^2+\alpha uv)$を持つ相対周囲はAS正則五次元で、
$W_{h,\alpha}=\Omega_h+\alpha uvx+u^3+v^3$は中心正則元となる。
中心性は$3D-h^2T+9mx$の三交換子の重み4の有限証明書で確認できる。
91語と84関係倍からなる行列のrankは69で、$\mathbb Q(h)$上の剰余は零となる。
使用したpivotの例外は$h(h^6-1)=0$だけである。
$h=0$は可換で、$h^6=1$では$h=1$での恒等式から
$(h,y)\mapsto(\omega h,\omega y)$（$\omega^3=1$）と反対代数により従う。
従って中心性を一般の$h$だけから例外値へ無条件に延長してはいない。
さらに$B_{h,\alpha}$は全ての$h,\alpha$で5-Calabi--Yauである。
実際、前節と同じKoszul双対を$\mathbb Q(h,\alpha)$上で計算すると、
次元列は$(1,5,10,10,5,1,0)$、Frobenius対合の行列式は
$(1,h^{-2},-h^{-6},-h^{-6},h^{-2},1)$となり、
次数1と次数4の積のtraceは可換である。
\texttt{verify\_central\_quotient\_koszul.py --general-h}は
全消去pivotの分子・分母の因子が$h(h^3+1)$を割ることも検証する。
従って$h\ne0$, $h^3\ne-1$では全ての$\alpha$で
Nakayama自己同型が恒等であり、一般点だけの計算ではない。
$h=0$は五変数多項式環、$h^3=-1$は$h=1$の場合の反対代数と
上記の三乗根による変数変換なので、例外値でも同じ結論を得る。
ここから$h^6\ne-27$とし、
$A_{\alpha,h}=B_{h,\alpha}/(W_{h,\alpha})$、
$C_{\alpha,h}=A_{\alpha,h}/(v-u)$とする。

\begin{corollary}[超越的な混合係数の具体例]
$\alpha$が$\mathbb Q(h)$上超越的なら
\[
 \operatorname{gldim}\operatorname{qgr}C_{\alpha,h}=2,\qquad
 \operatorname{gldim}\operatorname{qgr}A_{\alpha,h}=3.
\]
特に$(h,\alpha)=(1,\pi)$は非零混合項を持つ正則例である。
\end{corollary}
\begin{proof}
$\alpha\ne0$なので$x=\sqrt\alpha X,y=\sqrt\alpha Y,z=\sqrt\alpha Z$と規格化すると、
曲面の$u$-chartは
$R_h/(G_h+\beta)$、$G_h=\Omega_h+X$、$\beta=2\alpha^{-3/2}$となる。
$\beta$は$\mathbb Q(h)$上超越的で、$R_h$はPBWにより有限大域次元を持つ。
補題\ref{lem:global-transcendental}と命題\ref{prop:gluing-one-chart}を使う。
境界は$S_h/(\Omega_h)$の滑らかな楕円曲線で、
AS--Gorenstein次元3、index 1の双対性から曲面の正確な次元2を得る。
三様体は$u\partial_u-v\partial_v$と三切断$v=\omega u$を使う
命題\ref{prop:gluing-derivation}でperfectとなり、次元3が従う。
\end{proof}

\section{有限位数境界と非空Zariski開な正則範囲}
\label{sec:quotient-open-regularity}
以下の存在定理は、指定した一混合族の有効な必要十分判別式とは異なる。
\begin{lemma}[被約基底上の中心の持上げ]
\label{lem:global-center-lifting}
$H=B(E,\mathcal L,\tau)$は標数零の滑らかな楕円曲線のTHCRで、
$\mathcal L$はample、$\tau$は有限位数のtranslationとする。
$R$を被約有限型可換代数とし、filtered $R$-代数$U$が
$\operatorname{gr}U=H\otimes R$、各$F_dU$は有限自由で、
閉点への基底変更と両立するとする。
このとき各$z\in Z(H)_d$は$Z(U)\cap F_dU$へ持ち上がり、
$U$は有限型可換中心部分代数上有限となる。
\end{lemma}
\begin{proof}
斉次生成元の持上げ$u_i$を取り、有限自由加群間の交換子写像
$M_d:F_dU\to\bigoplus_iF_{d+\deg u_i}U$を作る。
Rains\cite[Theorem~2.13]{RainsElliptic}により各閉点$s$で
$\operatorname{gr}Z(U_s)=Z(H)$なので、$M_d$のrankは一定である。
被約性により高すぎるminorは恒等的に零となり、
最大rankのminorで被覆すれば$\ker M_d$は有限射影で、形成は基底変更と可換する。
最高次symbolの写像$\ker M_d\to Z(H)_d\otimes R$は
各閉点で全射、従ってNakayamaにより全射である。
右辺が自由なので、各斉次中心元は大域的に持ち上がる。
$H$は有限位数translationのnorm line bundleによる中心上有限である。
中心の有限生成元と中心上の有限加群生成元を持ち上げ、次数を下げる帰納法で
$U$も持ち上げ中心部分代数上有限と分かる。
ここではRainsの体上の定理を非被約基底へ拡張していない。
\end{proof}

\begin{theorem}[有限位数境界を持つ固定ambientの開な正則範囲]
\label{thm:global-open}
定理\ref{thm:global-very-general}の良い固定ambientを取り、さらに
$H=S/(\gamma)=B(E,\mathcal L,\tau)$のtranslationが有限位数とする。
ある非零多項式$P(a,b,c,d)$が存在して、$P\ne0$上では
中心三次商のtails大域次元が3となる。
周囲のPBW/Koszul解消とnoetherian AS条件が基底変更と両立する構成では、
この十分条件は任意の代数閉拡大体上でも成立する。
特に周囲の表示が有理数係数なら、有理数四係数の正則な商がZariski稠密に存在する。
\end{theorem}
\begin{proof}
$D_u=B[u^{-1}]_0$、$t=v/u$、$G_u=G/u^3$とする。
全係数を動かす$u$-chartの全代数は
\[
 \mathcal U=D_u[a,b,c,d]/(G_u+a+bt+ct^2+dt^3)
 \simeq D_u[b,c,d]
\]
なので左右noetherianかつ有限大域次元である。
$R=\CC[t,a,b,c,d]$をfiltration次数零とすると
$\operatorname{gr}\mathcal U=H\otimes R$で、各filtered pieceは有限自由。
前補題から$\mathcal U$は有限型可換中心部分代数$C_u$上有限である。
この$C_u$は係数環$R_0=\CC[a,b,c,d]$上も有限型である。
ここで\cite[Proposition~2.1, proof steps (6)--(7)]{RSVBertini}を
$(R,S,A)=(R_0,C_u,\mathcal U)$へ適用する。
その証明はある非零$P_u\in R_0$を逆にした上で、$\mathcal U$が$R_0$上平坦で、
相対対角が有限射影解消を持つことを与える。
従って全ての幾何学的fiberで有限大域次元が成立し、体拡大も許される。
$v$-chartでも$d$を消去して$P_v$を得る。
$P=P_uP_v$とし、二chartの貼合せとSerre双対性で結論する。
$\mathbb Q^4$は複素四次元affine空間でもZariski稠密なので、
非空開集合には有理数係数点が稠密に存在する。
\end{proof}
\begin{remark}
これは正則locus全体の開性を証明したものではなく、$P$の零点全てが悪いとも述べない。
$P$の係数・次数と一般混合族の有効判別式は未確定である。
$h=1$では境界translationの位数は6である
\cite[\S8.2]{WWYSklyanin}。
一混合族ではさらに本稿の正確な判別式が有限例外と全有理数点の正則性を与え、
この一次元部分族の正則locus自体はZariski開となる。
一般の無限位数境界ではvery-generalな十分条件までであり、全正則locusの開性は未解決。
\end{remark}

\begin{lemma}[位数6の具体的な楕円曲線計算]
$E_1=V(x^3+y^3+z^3+xyz)$で$O=[1:-1:0]$を原点とすると、
$p=[1:-1:-1]$は位数6であり、$2p=[1:0:-1]$である。
\end{lemma}
\begin{proof}
$E_1$は滑らかで、$O$は変曲点である。
$p$の接線$2x+y+z=0$へ代入すると
$7(y+z)(y-z)^2/8$となり、第三交点は$[0:1:-1]$である。
この点と$O,[1:0:-1]$は共線なので、接線の第三交点を反転すると表示した$2p$を得る。
$2p$は$O$と異なる変曲点なので位数3、$p$自体はHessianが非零で変曲点ではない。
従って$p$の位数は6である。各代入は現行の厳密検算に含まれる。
\end{proof}


% 旧中心商特異性論文VI 第15.3・16節の現行TeX回復
% 有限中心・Azumaya・中心chart gluingは現行証明を参照し、PBW/CY等の基礎入力を明示する。
\section{一次元表現の非被約性と正確なtails正則判別}
\label{sec:quotient-exact-tails}
この節は旧論文VIの第15.3節・第16節の証明を現在有効な順序で整理する。
結論は$h=1$の指定した一混合項族に限る。
全21混合項の一般判別式や任意の量子パラメータ$h$の判別式を
主張しない。有限中心・Azumaya locusは第\ref{sec:quotient-center-azumaya}節、
中心chartの貼合せは第\ref{sec:quotient-chart-gluing}節の証明を用いる。
PBW、周囲のAS正則性、$R_1$のCalabi--Yau性は
第\ref{sec:quotient-pbw-cy}節で証明した。
斉次Sklyaninの楕円境界の記述は既知入力として区別する。

$R_1=\CC\langle X,Y,Z\rangle/
([X,Y]-Z^2,[Y,Z]-(X^2+1),[Z,X]-Y^2)$、
$G_1=\Omega_1(X,Y,Z)+X$とし、
$D_\beta=R_1/(G_1+\beta)$と置く。

\begin{proposition}[一次元表現を検出する中心値]
\label{prop:tails-one-dimensional}
$D_\beta$に一次元表現があるための必要十分条件は
$9\beta^2+4=0$である。
$s^2=-1$、$\beta=-2s/3$なら、
唯一の一次元表現の局所スキームは
\[
 \operatorname{Rep}_1(D_{-2s/3})
 \simeq\operatorname{Spec}\frac{\CC[y,z]}{(y^2,yz,z^2)}
\]
であり、長さ3、接空間次元2である。
全$\beta$にわたる一次元表現のスキーム論的像の
定義イデアルは$((9\beta^2+4)^2)$である。
\end{proposition}
\begin{proof}
可換化を取ると
$T=\CC[x,y,z]/(x^2+1,y^2,z^2)$となり、
$G_1$の像は$\frac23x(1-yz)$である。
$x^2+1$の二根$s=\pm i$の各因子で
$G_1+\beta=\beta+\frac23s-\frac23syz$となる。
前半はこの式の単元性を調べれば得られる。
さらに$\beta=-2s/3$では$yz=0$が追加され、
所望の局所環になる。
全因子にわたる中心写像は
$\beta=-\frac23x(1-yz)$である。
$x^2=-1$および$(yz)^2=0$から
$9\beta^2+4=8yz$であり、
各因子で$yz\ne0$なので像の核は
$((9\beta^2+4)^2)$となる。
\end{proof}

\begin{lemma}[有限中心を持つ中心正則商のCY双対性]
\label{lem:tails-central-cy}
$R$を$\CC$上の$d$次元Calabi--Yau代数、
$w\in Z(R)$を左右正則元とし$D=R/(w)$とする。
$D$のhomological smoothnessを仮定しなくても、
双対化複体には
\[
 \mathbf R\operatorname{Hom}_{D^e}(D,D^e)
 \simeq D[-(d-1)]
\]
という同型がある。
さらに$D$が有限生成可換中心部分代数上有限で
有限大域次元を持てば、
$D$はhomologically smoothで$(d-1)$-Calabi--Yauである。
\end{lemma}
\begin{proof}
$R^e=R\otimes R^{\mathrm{op}}$内の中心正則列
$w\otimes1,1\otimes w$のKoszul分解から
$\mathbf R\operatorname{Hom}_{R^e}(D^e,R^e)
 \simeq D^e[-2]$を得る。
一方、$0\to R\xrightarrow{w}R\to D\to0$と
$\mathbf R\operatorname{Hom}_{R^e}(R,R^e)\simeq R[-d]$
から
$\mathbf R\operatorname{Hom}_{R^e}(D,R^e)\simeq D[-d-1]$。
導来随伴で二つを比較すれば最初の表示が従う。
有限型可換中心部分代数$C$上有限なら、$D$と$D^{\mathrm{op}}$は
affine PIかつ左右noetherianで、$D^e$も$C\otimes C$上有限だからnoetherian。
有限大域次元を仮定した場合、
\cite[Lemmas~3.4, 3.5]{RSVBertini}から
$D^{\mathrm{op}}$と$D^e$も有限大域次元を持つ。
ここで基礎体$\CC$の代数閉性とaffine PI性を使った。
従って有限生成$D^e$加群$D$はperfectでsmoothとなる。
非可換中心商が無条件にsmoothになるとは主張していない。
\end{proof}

\begin{theorem}[二つの真に悪いaffine中心値]
\label{thm:tails-bad-affine}
命題\ref{prop:quotient-constant-cy}と
第\ref{sec:quotient-center-azumaya}節により、
$R_1$は三次元Calabi--Yauで中心上有限、$G_1+\beta$は中心正則元である。
このとき
\[
 \beta=\pm\frac{2i}{3}
 \quad\Longrightarrow\quad \operatorname{gldim}D_\beta=\infty.
\]
\end{theorem}
\begin{proof}
もし大域次元が有限なら、
補題\ref{lem:tails-central-cy}から$D_\beta$は
二次元Calabi--Yauとなる。

Bocklandt--Galluzzi--\allowbreak Vaccarinoの単純表現の局所構造定理
\cite[Theorems 6.3, 6.6 and Corollary 6.8]{BGV}
によれば、その一次元単純表現の点で
$\operatorname{Rep}_1(D_\beta)$は滑らかである必要がある。
しかし命題\ref{prop:tails-one-dimensional}の局所環は
零次元・接空間次元2であり非正則だから矛盾する。
\end{proof}

\begin{lemma}[有限射影行列因子化によるnonperfectness検出]
\label{lem:tails-finite-mf}
$R$を左右noetherian、$\operatorname{gldim}R\le d<\infty$、
$w\in Z(R)$を正則元とし、$D=R/(w)$と置く。
有限生成$D$加群の$(d-1)$次syzygyは
有限生成射影$R$加群の行列因子化のcokernelとして得られる。
さらにそのcokernelが有限射影次元なら
当該行列因子化はhomotopy縮約可能である。
従って$\operatorname{gldim}D=\infty$なら
非縮約な有限射影行列因子化が存在する。
\end{lemma}
\begin{proof}
$\operatorname{pd}_R D=1$なので、有限自由$D$分解を
$(d-1)$回ずらしたsyzygy $N$は
$R$上の長さ1の有限射影分解
$0\to P_1\xrightarrow{\phi}P_0\to N\to0$を持つ。
$wN=0$から$\phi\psi=w$を満たす$\psi$を得て
単射性から$\psi\phi=w$となる。
$D$上で有限射影次元なら補題\ref{lem:gluing-contraction}で縮約できる。
全ての有限生成加群の射影次元が有限なら同じ構成で
上界$d-1$を一様に与え、最後の主張が従う。
\end{proof}

\begin{lemma}[中心値の二重分岐はnonperfectnessを保存]
\label{lem:tails-double-zero}
$R$は左右noetherianで有限大域次元を持ち、
$G+\beta_0\in Z(R)$は正則で
$\operatorname{gldim}R/(G+\beta_0)=\infty$とする。
$s_0\in\CC^\times$、$b(s)\in\CC[s,s^{-1}]$が
\[
 b(s)-\beta_0=a(s)c(s),\quad
 a(s_0)=c(s_0)=0
\]
を満たし、$G+b(s)$は$R[s,s^{-1}]$上正則とする。
このとき
$T=R[s,s^{-1}]/(G+b(s))$は有限生成非perfect加群を持つ。
\end{lemma}
\begin{proof}
補題\ref{lem:tails-finite-mf}より
$G+\beta_0$の非縮約な有限射影行列因子化
$(P,d_0)$を取る。
これを$R[s,s^{-1}]$に拡大し、
中心的階数1の因子化$(a,c)$とsuper tensorを取る。
全微分$d$は$d^2=G+\beta_0+ac=G+b(s)$を満たす。
そのcokernelがperfectなら前補題の縮約の議論により
$d$は縮約できる。
$s=s_0$へ特殊化すると$(a,c)$は零微分となり、
$d_0\oplus d_0[1]$が縮約する。
直和因子に射影すれば$d_0$も縮約して矛盾。
\end{proof}

\begin{lemma}[tailsから中心chart、有限\'etale基底変更へ]
\label{lem:tails-chart-perfect}
$A$をnoetherian connected graded代数とし、
$u\in Z(A)_1$を正則元とする。
$\operatorname{gldim}\operatorname{qgr}A<\infty$なら
$Q_0=A[u^{-1}]_0$の任意の有限生成加群はperfectである。
この性質は中心局所化、
および有限\'etale可換中心拡大$Q\to T$にも伝わる。
\end{lemma}
\begin{proof}
$Q_0$の有限生成加群$M$の有限個の生成元から、
$M[u,u^{-1}]$内に有限生成graded $A$-latticeを取れる。
$\operatorname{qgr}A$の有限大域次元とtwistsの生成により、
その有界導来圏は$\pi A(m)$で厚く生成される。
中心局所化後、twistsは$Q_0$上自由だから$M$はperfect。
中心局所化には有限生成latticeが存在しperfect性を保つ。
有限\'etale拡大では$M$は
$M\otimes_QT$の$T$-直和因子であり、
separability idempotentによりperfect性が降下する。
\end{proof}

\begin{proposition}[曲面の判別式：中心chartからの導出]
\label{prop:tails-surface-discriminant}
$C_\alpha=A_\alpha/(v-u)$とし、
$B_\alpha,A_\alpha$の定義を
第\ref{sec:quotient-exact-tails}節の後半の定理に置く。
第\ref{sec:quotient-pbw-cy}節のPBW・AS正則性・CY性と、
境界$\operatorname{qgr}(C_\alpha/(u))\simeq\operatorname{coh}E_1$
の既知の記述（$E_1$は滑らかな楕円曲線）を用いる。
貼合せ自体を追加仮定とはしない。
すると
\[
 \operatorname{gldim}\operatorname{qgr}C_\alpha=
 \begin{cases}2,&(\alpha^3+9)(28\alpha^3+243)\ne0,\\
 \infty,&(\alpha^3+9)(28\alpha^3+243)=0.\end{cases}
\]
\end{proposition}
\begin{proof}
$\alpha\ne0$では中心chartは
$R_1/(G_1+\beta)$、$\beta^2\alpha^3=4$に同型。
中心曲面判別式とAzumaya性から
$\gamma^2\notin\{0,-1,-28/27\}$では
$\operatorname{gldim}R_1/(g-\gamma)=2$である
（$g=3G_1/2$、$\gamma=-3\beta/2$）。
$\gamma^2=-28/27$では前節の中心Azumaya点における
特異曲面から無限大域次元、
$\beta=\pm2i/3$では
定理\ref{thm:tails-bad-affine}から無限大域次元。
後二者は補題\ref{lem:tails-chart-perfect}の逆向き対偶により
tailsにも障害となる。
よい値では楕円境界は有限大域次元を持ち、
そのcoherent対象は構造対象のtwistsで厚く生成される。
実際、twistsの和による解消を大域次元より一つ長く切り、
最後のYoneda拡大の消滅から各対象を有限複体の直和因子として得る。
命題\ref{prop:gluing-one-chart}で全coherent対象のperfect性を得て、
AS--Gorenstein次元3、index 1の双対化複体に
補題\ref{lem:gluing-serre-dimension}を適用するとtails大域次元は正確に2となる。
$\alpha=0$では$u$-chartが$S_1/(\Omega_1+2)$となる。
Euler導分$\delta(\Omega_1)=3\Omega_1$に
補題\ref{lem:gluing-moving-central-value}を適用すると、
$P(-2)=-6\ne0$からchartの大域次元は高々2となる。
同じ命題とSerre双対性が適用できる。
\end{proof}

\begin{theorem}[一混合項族の三次三様体と曲面の正確な判別式]
\label{thm:quotient-exact-tails}
\[
 B_\alpha=
 \frac{\CC[u,v]\langle x,y,z\rangle}
 {([x,y]-z^2,\ [y,z]-(x^2+\alpha uv),\
   [z,x]-y^2)},\qquad
 A_\alpha=B_\alpha/(W_\alpha),
\]
\[
 W_\alpha=\frac{x^3+y^3+z^3}{3}
 -\frac{xyz+xzy+yxz+yzx+zxy+zyx}{9}
 +\alpha uvx+u^3+v^3,\quad C_\alpha=A_\alpha/(v-u)
\]
とする。系\ref{cor:quotient-ambient-pbw}と
命題\ref{prop:quotient-constant-cy}により、
$B_\alpha$はKoszul AS正則（次元5、index 5）、
$W_\alpha$は中心正則、$A_\alpha$は$\CC[u,v]$上自由で、
$R_1$は三次元Calabi--Yauである。
斉次Sklyaninの既知の記述
$\operatorname{qgr}(S_1/(\Omega_1))\simeq\operatorname{coh}E_1$
（$E_1$は滑らかな楕円曲線）を用いる。
最後の境界の代数的表示は
\cite[Lemma~2.19]{WWYSklyanin}のtwisted homogeneous coordinate ringである。
中心chartの貼合せは第\ref{sec:quotient-chart-gluing}節で証明済みである。
$\Delta(\alpha)=(\alpha^3+9)(28\alpha^3+243)$と置くと
\[
 \bigl(\operatorname{gldim}\operatorname{qgr}C_\alpha,\,
 \operatorname{gldim}\operatorname{qgr}A_\alpha\bigr)
 =
 \begin{cases}(2,3),&\Delta(\alpha)\ne0,\\
 (\infty,\infty),&\Delta(\alpha)=0.
 \end{cases}
\]
従って全ての$\alpha\in\mathbb Q$は正則であり、
$\alpha=1$は有理係数の具体的な本質混合例を与える。
\end{theorem}
\begin{proof}
曲面は前命題による。
三次三様体のよい値では、
中心導分$u\partial_u-v\partial_v$の
$\delta(W_\alpha)=3(u^3-v^3)$と、
中心切断$v=\omega u$（$\omega^3=1$）を用いる。
この導分は$uv$を固定するため、関係式を保つ真の導分である。
$\delta(W_\alpha)=-3\prod_{\omega^3=1}(v-\omega u)$であり、
各因子は$\CC[u,v]$上の自由性から$A_\alpha$上正則。
自己同型$(u,v)\mapsto(\zeta u,\zeta^{-1}v)$（$\zeta^3=1$）は
周囲の関係式と$W_\alpha$を保ち、三つの切断を巡回的に交換する。
各切断は前命題の$C_\alpha$と同型なので、
命題\ref{prop:gluing-derivation}から全coherent対象のperfect性を得る。
$A_\alpha$のAS--Gorenstein次元4、index 2を
補題\ref{lem:gluing-serre-dimension}に適用すれば、
$\operatorname{gldim}\operatorname{qgr}A_\alpha=3$となる。

悪い値では$\alpha\ne0$である。
$u,v$を逆にした次数零chartで$t=v/u$とし、
中心に$s^2=\alpha t$を添加する。
$s,2s$が可逆なのでこれは有限\'etale拡大となり、
$x=sX,y=sY,z=sZ$とすると
\[
 T\simeq R_1[s,s^{-1}]/(G_1+b_\alpha(s)),
 \qquad b_\alpha(s)=s^{-3}+\frac{s^3}{\alpha^3}.
\]
$s_0^2=\alpha$とし$\beta_0=2/s_0^3$と置く。
$\beta_0^2\alpha^3=4$であり、
$\Delta(\alpha)=0$なら
$\beta_0^2=-4/9$または$-112/243$。
前者は定理\ref{thm:tails-bad-affine}、
後者は前節のAzumaya局所判定から
$\operatorname{gldim}R_1/(G_1+\beta_0)=\infty$。
一方
\[
 b_\alpha(s)-\beta_0
 =\frac{(s^3-s_0^3)^2}{s_0^6s^3}.
\]
ここで$a(s)=s^3-s_0^3$、
$c(s)=a(s)/(s_0^6s^3)$はともに$s=s_0$で零。
補題\ref{lem:tails-double-zero}により$T$に
有限生成非perfect加群が存在する。
もし$\operatorname{gldim}\operatorname{qgr}A_\alpha<\infty$なら
補題\ref{lem:tails-chart-perfect}により$T$の全有限生成加群はperfectで、
矛盾する。これで必要性を得る。
最後に$\alpha\in\mathbb Q$なら
$\alpha^3=-9,-243/28$は三進付値により不可能。
\end{proof}

\begin{remark}[量子的判別式と古典・点加群の違い]
$\alpha^3=-243/28$では古典三次三様体と
量子的楕円境界は滑らかなまま、tailsのみ非正則となる。
さらに非零$\alpha$で$\alpha^3\ne-9$の範囲では
正則例と新しい非正則例の定常点加群スキームは
六点を台とする長さ4の局所環
$\CC[Y,Z]/(Y^2,Z^2)$を持つ。
これは埋込を含む点スキーム全体の同型を主張しない。
\end{remark}


\begin{corollary}[全affine中心値の必要十分判別式]
\label{cor:tails-all-affine}
全ての$\beta\in\CC$について
\[
 \operatorname{gldim}R_1/(G_1+\beta)
 =\begin{cases}
 2,&(9\beta^2+4)(243\beta^2+112)\ne0,\\
 \infty,&(9\beta^2+4)(243\beta^2+112)=0.
 \end{cases}
\]
\end{corollary}
\begin{proof}
$\beta=0$は系\ref{cor:zero}による。
それ以外は$g=-3\beta/2$、$H=9\beta^2/4$として
系\ref{cor:quotient-azumaya-fibers}の滑らかな中心ファイバーと特異Azumaya点の判定を使う。
$9\beta^2+4=0$の場合は定理\ref{thm:tails-bad-affine}による。
\end{proof}


\section{証明の入力と残る照合対象}
全中心の関係式は有理数上の有限計算で確認した。
全中心生成には既知の斉次Sklyanin中心定理を用い、
Azumaya locusにはPBW性、斉次Sklyaninの諸性質と
一次文献\cite{WWYSklyanin,LevasseurNoncomm,BrownMacleod}の
明記した定理を適用した。
Rees退化、中心の忠実平坦降下、15点のHessian、
三つのホモロジー的特異点の行列node模型、境界とchartのGorenstein性、
中心chartおよび導分によるperfect性の貼合せは現行の証明に含まれる。

相対PBW性、周囲のAS正則性、$R_1$の3-CY性と周囲のNakayama twistの恒等性は
第\ref{sec:quotient-pbw-cy}節に証明・有限証明書を収録した。
斉次Sklyaninの標準的性質と楕円境界の記述、一般のPBW・分類・双対性定理は既知入力である。
非Azumayaな四点でのrank-one cornerとmodifying generatorは、
六次分裂元と\cite[Theorem~5.1(1)]{IW}の仮定を照合して導いた。
小単純表現の行列と全表現間の一次Ext表は再現コードを持ち、
命題\ref{prop:quotient-full-ext1}にはExt表の構造的証明も与えた。
分離型三次式の必要十分条件、一般周囲での超越的中心値、
有限位数境界での正則開集合も、第\ref{sec:quotient-global-regularity}節以下に
仮定と証明を収録した。Cayleyコンパクト化と古典判別式も独立に計算した。
原資料全22節の対応と置換した弱い結果は\texttt{PDF\_SOURCES.md}に記す。
既知の一般定理の具体化と厳密計算を区別し、
具体式・構成の初出について網羅的な新規性を主張しない。
一般の量子パラメータでの有効判別式、Brauer類の大域的表示、
標識なし復元の存在分類はなお研究課題である。

\begin{thebibliography}{99}
\bibitem{ATVRegular}
M.~Artin, J.~Tate and M.~Van den Bergh,
\emph{Some algebras associated to automorphisms of elliptic curves},
The Grothendieck Festschrift, Vol.~I, Progr. Math. 86 (1990), 33--85.
\bibitem{ISSklyanin}
N.~Iyudu and S.~Shkarin, \emph{Sklyanin algebras and a cubic root of 1},
MPIM Preprint 2017 (49), Theorem~1.1;
\url{https://archive.mpim-bonn.mpg.de/997/1/preprint_2017_49.pdf}.
\bibitem{IW} O.~Iyama and M.~Wemyss,
\emph{Reduction of triangulated categories and maximal modification
algebras for cAn singularities},
J.\ Reine Angew.\ Math.\ \textbf{738} (2018), 149--202,
\href{https://doi.org/10.1515/crelle-2015-0031}{doi:10.1515/crelle-2015-0031};
\url{https://arxiv.org/abs/1304.5259}.
Theorem~5.1(1)（$f=(r-\gamma)(r+\gamma)$の二素因子とflag分類を適用）。
\bibitem{ChenRelativeSingcat}
X.-W. Chen, \emph{Relative singularity categories and
Gorenstein-projective modules}, Math. Nachr. 284 (2011), 199--212,
\url{https://arxiv.org/abs/0709.1762};
Theorem~3.8（stable categoryと特異圏の比較）.
\bibitem{OrlovSingcat}
D. Orlov, \emph{Triangulated categories of singularities and
D-branes in Landau--Ginzburg models},
Proc. Steklov Inst. Math. 246 (2004), 227--248,
\url{https://arxiv.org/abs/math/0302304}; Theorem~2.1.
\bibitem{SeidelThomas}
P. Seidel and R. P. Thomas,
\emph{Braid group actions on derived categories of coherent sheaves},
Duke Math. J. 108 (2001), 37--108,
\url{https://arxiv.org/abs/math/0001043};
Definition~2.14, Proposition~2.10, Lemma~2.15.
\bibitem{YZSerre}
A. Yekutieli and J. J. Zhang,
\emph{Serre duality for noncommutative projective schemes},
Proc. Amer. Math. Soc. 125 (1997), 697--707;
\url{https://doi.org/10.1090/S0002-9939-97-03782-9}.
\bibitem{RainsElliptic}
E.~M.~Rains, \emph{Filtered deformations of elliptic algebras},
\url{https://arxiv.org/abs/2107.13540}, v2 (2022), Theorem~2.13 and its preceding discussion.
\bibitem{EGDelPezzo}
P.~Etingof and V.~Ginzburg,
\emph{Noncommutative del Pezzo surfaces and Calabi--Yau algebras},
J. Eur. Math. Soc. 12 (2010), 1371--1416;
\url{https://arxiv.org/abs/0709.3593}, Theorems~3.3.2, 3.4.4, 3.4.5.
\bibitem{VdBOverview}
M.~Van den Bergh, \emph{Non-commutative crepant resolutions, an overview},
\url{https://arxiv.org/abs/2207.09703}, Definitions~2.1 and 2.2.
\bibitem{DKSResultant}
C.~D'Andrea, T.~Krick and A.~Szanto,
\emph{Multivariate subresultants in roots}, J. Algebra 302 (2006), 16--36;
\url{https://arxiv.org/abs/math/0501281}, equation~(2).
\bibitem{RRZNakayama}
M.~Reyes, D.~Rogalski and J.~J.~Zhang,
\emph{Skew Calabi--Yau algebras and homological identities},
Adv. Math. 264 (2014), 308--354;
\url{https://arxiv.org/abs/1302.0437},
Lemmas~1.2, 1.5 and proof of Theorem~6.3.
\bibitem{BGPBW}
A.~Braverman and D.~Gaitsgory,
\emph{The Poincar\'e--Birkhoff--Witt theorem for quadratic algebras of Koszul type},
J. Algebra 181 (1996), 315--328;
\url{https://arxiv.org/abs/hep-th/9411113}, Main Theorem~0.5 and Lemma~3.3.
\bibitem{RSVBertini}
J.~V.~Rennemo, E.~Segal and M.~Van den Bergh,
\emph{A non-commutative Bertini theorem},
\url{https://arxiv.org/abs/1705.01366}, v2 (2020), Lemmas~3.4 and 3.5.
\bibitem{CCKM}
T.~Cassidy, A.~Conner, E.~Kirkman and W.~F.~Moore,
\emph{Periodic free resolutions from twisted matrix factorizations},
J. Algebra 455 (2016), 137--163;
\url{https://arxiv.org/abs/1307.8415}, Theorems~4.7 and 5.5.
\bibitem{ChenGeneralizedSerre}
X.-W.~Chen, \emph{Generalized Serre duality},
\url{https://arxiv.org/abs/math/0610258}, Theorem~3.10.
\bibitem{WWYSklyanin}
C.~Walton, X.~Wang and M.~Yakimov,
\emph{Poisson geometry of PI 3-dimensional Sklyanin algebras},
\url{https://arxiv.org/abs/1704.04975};
Proposition~2.31, Theorem~1.3, Lemma~5.4 and \S8.2.
\bibitem{LevasseurNoncomm}
T.~Levasseur,
\emph{Some properties of non-commutative regular graded rings},
Glasgow Math. J. 34 (1992), 277--300,
\url{https://doi.org/10.1017/S0017089500008843};
Proposition~3.2(ii) and \S5.
\bibitem{BrownMacleod}
K.~A.~Brown and M.~J.~Macleod,
\emph{The Cohen Macaulay property for noncommutative rings},
Algebr. Represent. Theory 20 (2017), 1433--1465,
\url{https://doi.org/10.1007/s10468-017-9694-z};
Theorems~3.13 and 4.8.
\bibitem{BGV}
R.~Bocklandt, F.~Galluzzi and F.~Vaccarino,
\emph{The Nori--Hilbert scheme is not smooth for 2-Calabi--Yau algebras},
J. Noncommut. Geom. \textbf{10} (2016), 745--774,
\url{https://doi.org/10.4171/JNCG/247};
Theorems~6.3, 6.6, 6.7 and Corollary~6.8.
\end{thebibliography}
\end{document}
